\begingroup
\tcbset{breakable/.default=false,after={\par\nobreak\smallskip}}
\RenewDocumentEnvironment{xesolution}{+b}{%
  \par\nobreak\smallskip\nobreak\begingroup\small
  \clubpenalties 2 10000 10000\displaywidowpenalties 2 10000 10000\noindent
  {\bfseries\color{DeepBlue}解答}\quad #1\par\endgroup\smallskip
}{}
\begingroup
\renewcommand{\Needspace}[1]{\par\begingroup
\dimen0=#1\relax\vskip0pt plus\dimen0\penalty-100
\vskip0pt plus-\dimen0\vskip\dimen0\penalty9999\vskip-\dimen0
\vskip0pt\endgroup}
\tcbset{breakable/.default=false}
\chapter{恒定电流与复杂电路}

{
\graphicspath{{/Users/dongxingxi/Documents/物理竞赛一轮资料/sample/xes-original/thirty-sixth-section/}{/Users/dongxingxi/Documents/物理竞赛一轮资料/sample/xes-original/thirty-sixth-section/assets/}}
\section{电流密度与微观输运}

% All labels are transparent. Rectangles in circuit figures represent actual elements,
% not label backgrounds.
\tikzset{
  xeslabel/.style={font=\small,text=Ink,fill=none,inner sep=.55pt},
  xessmall/.style={font=\footnotesize,text=TextGray,fill=none,inner sep=.4pt},
  xestiny/.style={font=\scriptsize,text=TextGray,fill=none,inner sep=.3pt},
  xesline/.style={draw=Ink,line width=.85pt,line cap=round,line join=round},
  xesblue/.style={draw=DeepBlue,line width=1pt,line cap=round,line join=round},
  xesarrow/.style={-{Latex[length=2.1mm,width=1.3mm]},draw=DeepBlue,line width=.9pt},
  grayarrow/.style={-{Latex[length=1.7mm,width=1.1mm]},draw=TextGray,line width=.65pt},
  circuit/.style={draw=Ink,line width=.9pt,line cap=round,line join=round},
  element/.style={draw=Ink,fill=PaleBlue,minimum width=.66cm,minimum height=.30cm,
    inner sep=0pt,line width=.75pt},
  nodepoint/.style={circle,fill=DeepBlue,draw=none,inner sep=1.8pt},
  driftpath/.style={draw=DeepBlue!75,line width=.65pt}
}


\newcommand{\ivCharacteristicFigure}{%
\begin{center}
\begin{tikzpicture}[x=1cm,y=1cm]
  \begin{scope}[shift={(-3.2,0)}]
    \draw[circuit] (-2.15,0)--(-.75,0);
    \draw[circuit,fill=DeepBlue!7] (-.75,-.42) rectangle (.75,.42);
    \draw[circuit] (.75,0)--(2.15,0);
    \fill[white,draw=Ink,line width=.8pt] (-2.15,0) circle (2pt);
    \fill[white,draw=Ink,line width=.8pt] (2.15,0) circle (2pt);
    \draw[xesarrow] (-1.95,.62)--(-.95,.62);
    \node[xeslabel,above] at (-1.45,.65) {$I$};
    \node[xeslabel,below] at (0,-.48) {$V$};
    \node[xeslabel,below] at (-2.15,-.12) {$+$};
    \node[xeslabel,below] at (2.15,-.12) {$-$};
    \node[xessmall] at (0,-1.15) {二端元件与人为规定的测试方向};
  \end{scope}
  \begin{scope}[shift={(3.25,0)}]
    \draw[xesline,-{Latex[length=1.8mm]}] (-2.05,0)--(2.05,0) node[xeslabel,right] {$V$};
    \draw[xesline,-{Latex[length=1.8mm]}] (0,-1.55)--(0,1.65) node[xeslabel,above] {$I$};
    \draw[xesblue] (-1.65,-1.22)..controls(-1.20,-1.12)and(-.92,-.72)..(-.62,-.50)
      ..controls(-.20,-.18)and(.38,.24)..(.73,.53)
      ..controls(1.05,.79)and(1.28,1.13)..(1.43,1.48);
    \node[xessmall] at (0,-1.90) {伏安关系 $I=I(V)$};
  \end{scope}
\end{tikzpicture}
\end{center}}

\newcommand{\ohmHistoryFigure}{%
\begin{center}
  \includegraphics[width=.88\linewidth]{ohm-portrait-and-stamp-built-in-v1.png}
  \par\smallskip
  \begin{minipage}[t]{.43\linewidth}\centering\small\color{TextGray}
    (a) 格奥尔格·西蒙·欧姆（1789--1854）的历史肖像。
  \end{minipage}\hfill
  \begin{minipage}[t]{.43\linewidth}\centering\small\color{TextGray}
    (b) 以电阻和 $U=RI$ 为主题的纪念邮票。
  \end{minipage}
  \figcap{欧姆肖像与电阻定律纪念邮票。}\label{ch21:fig:s36-5-2}
\end{center}}

\newcommand{\cylindricalResistorFigure}{%
\begin{center}
\begin{tikzpicture}[x=1cm,y=1cm]
  \fill[DeepBlue!8] (-2.4,-.55) rectangle (2.4,.55);
  \draw[xesblue] (-2.4,-.55)--(2.4,-.55) (2.4,.55)--(-2.4,.55);
  \draw[xesblue] (-2.4,0) ellipse (.26 and .55);
  \draw[xesblue] (2.4,.55) arc[start angle=90,end angle=-90,x radius=.26,y radius=.55];
  \draw[xesblue,dashed] (2.4,-.55) arc[start angle=-90,end angle=-270,x radius=.26,y radius=.55];
  \draw[<->,Gold!70!black,line width=.65pt] (-2.35,.95)--(2.35,.95);
  \node[xestiny,text=Gold!70!black] at (0,1.18) {$l$};
  \draw[xesarrow] (-3.15,0)--(-2.48,0);
  \node[xeslabel,above] at (-2.82,.06) {$I$};
  \node[xeslabel] at (0,0) {$\rho$};
  \draw[RuleGray,line width=.55pt,-{Latex[length=1.5mm]}] (-3.45,-.85)--(-2.58,-.32);
  \node[xessmall,anchor=east] at (-3.50,-.88) {截面积 $A$};
  \node[xeslabel,below] at (-2.40,-.65) {$+$};
  \node[xeslabel,below] at (2.40,-.65) {$-$};
\end{tikzpicture}
\end{center}}

\newcommand{\ohmicIVFigure}{%
\begin{center}
\begin{tikzpicture}[x=1cm,y=1cm]
  \draw[xesline,-{Latex[length=1.8mm]}] (-2.15,0)--(2.15,0) node[xeslabel,right] {$V$};
  \draw[xesline,-{Latex[length=1.8mm]}] (0,-1.45)--(0,1.65) node[xeslabel,above] {$I$};
  \draw[xesblue] (-1.65,-1.25)--(1.72,1.30);
  \draw[RuleGray,densely dashed] (1.10,0)--(1.10,.83)--(0,.83);
  \node[xessmall,anchor=west] at (1.18,.58) {斜率 $1/R$};
  \node[xestiny,below left] at (0,0) {$O$};
\end{tikzpicture}
\end{center}}

\newcommand{\drudeFigure}{%
\begin{center}
\begin{tikzpicture}[x=1cm,y=1cm]
  \fill[DeepBlue!4] (-4.2,-1.45) rectangle (4.2,1.45);
  \foreach \x/\y in {-3.6/1.0,-2.4/.95,-1.15/1.05,.2/.95,1.45/1.08,2.75/.92,3.70/1.10,
                        -3.1/.05,-1.85/.15,-.55/-.02,.85/.12,2.15/.0,3.35/.15,
                        -3.65/-.95,-2.35/-1.02,-1.0/-.90,.35/-1.05,1.65/-.92,2.95/-1.06,3.82/-.92}{
    \fill[Gold!78] (\x,\y) circle (3.8pt);\draw[Gold!60!black,line width=.45pt] (\x,\y) circle (3.8pt);}
  \draw[driftpath] (-4.0,-.58)..controls(-3.55,-.20)and(-3.30,.45)..(-3.05,.10)
    ..controls(-2.70,-.30)and(-2.25,-.42)..(-2.02,.02)
    ..controls(-1.75,.50)and(-1.18,.35)..(-.96,-.05)
    ..controls(-.65,-.62)and(-.15,-.42)..(.10,.02)
    ..controls(.38,.55)and(.95,.46)..(1.18,.05)
    ..controls(1.45,-.42)and(2.0,-.52)..(2.28,-.04)
    ..controls(2.55,.42)and(3.12,.30)..(3.95,.58);
  \foreach \x/\y in {-3.80/-.48,-3.08/.10,-2.03/.02,-.96/-.05,.10/.02,1.18/.05,2.28/-.04,3.78/.49}{
    \fill[DeepBlue] (\x,\y) circle (2pt);}
  \draw[xesarrow] (-3.45,-1.82)--(-1.55,-1.82);
  \node[xessmall,below] at (-2.50,-1.88) {电子平均漂移 $\boldsymbol v_d$};
  \draw[xesarrow] (3.45,-1.82)--(1.55,-1.82);
  \node[xessmall,below] at (2.50,-1.88) {电场与电流 $\boldsymbol E,\boldsymbol j$};
  \node[xessmall] at (0,1.82) {高速无规则运动上叠加很小的定向漂移；每次碰撞近似“重置”漂移速度};
\end{tikzpicture}
\end{center}}

\newcommand{\gasDischargeFigure}{%
\begin{center}
\begin{tikzpicture}[x=.82cm,y=.78cm]
  \draw[xesline,-{Latex[length=1.8mm]}] (0,0)--(10.1,0) node[xessmall,right] {电流 $I$（对数标度）};
  \draw[xesline,-{Latex[length=1.8mm]}] (0,0)--(0,5.55) node[xessmall,above] {电压 $U$};
  \foreach \x/\lab in {1/{10^{-15}},3.7/{10^{-10}},6.2/{10^{-5}},8.7/{10^{0}}}{
    \draw[RuleGray] (\x,0)--(\x,-.10);\node[xestiny,below] at (\x,-.12) {$\lab$};}
  \foreach \y/\lab in {1/{200},2/{400},3/{600},4/{800},5/{1000}}{
    \draw[RuleGray] (0,\y)--(-.10,\y);\node[xestiny,anchor=east] at (-.14,\y) {$\lab$};}
  \draw[RuleGray,densely dashed] (5.72,0)--(5.72,5.15) (8.55,0)--(8.55,5.15);
  \node[xessmall] at (2.80,5.20) {暗放电};
  \node[xessmall] at (7.15,5.20) {辉光放电};
  \node[xessmall] at (9.25,5.20) {电弧};
  \draw[xesblue] (.55,.42)..controls(.95,.42)and(1.05,2.65)..(1.34,2.95)
    ..controls(1.58,3.12)and(4.8,3.20)..(5.70,3.26)
    ..controls(6.18,3.22)and(6.62,2.20)..(7.23,1.64)
    ..controls(7.55,1.34)and(7.96,1.50)..(8.18,2.20)
    ..controls(8.38,2.90)and(8.43,4.72)..(8.55,5.05);
  \draw[Gold!85!black,line width=1.05pt] (8.55,5.05)..controls(8.67,2.40)and(8.72,.72)..(9.05,.62)
    ..controls(9.36,.55)and(9.72,.66)..(9.92,.74);
  \foreach \x/\y/\lab in {.58/.45/{A'},3.10/3.16/{C},5.68/3.26/{D},6.15/2.92/{E},6.85/1.95/{F},7.45/1.46/{G},8.12/2.22/{H},8.55/5.05/{I},9.08/.62/{J},9.85/.74/{K}}{
    \fill[Ink] (\x,\y) circle (1.7pt);\node[xestiny,above=2pt] at (\x,\y) {$\lab$};}
  \fill[Ink] (1.35,2.95) circle (1.7pt);\node[xestiny,below left=1pt] at (1.35,2.95) {$B'$};
  \fill[Ink] (1.65,3.08) circle (1.7pt);\node[xestiny,above right=1pt] at (1.65,3.08) {$B$};
\end{tikzpicture}
\end{center}}

\newcommand{\seriesParallelResistorsFigure}{%
\begin{center}
\begin{tikzpicture}[x=1cm,y=1cm]
  \begin{scope}[shift={(-3.55,0)}]
    \draw[circuit] (-2.05,0)--(-1.35,0);
    \draw[circuit,decorate,decoration={zigzag,segment length=2.8mm,amplitude=1.05mm}] (-1.35,0)--(-.10,0);
    \draw[circuit] (-.10,0)--(.35,0);
    \draw[circuit,decorate,decoration={zigzag,segment length=2.8mm,amplitude=1.05mm}] (.35,0)--(1.60,0);
    \draw[circuit] (1.60,0)--(2.05,0);
    \node[xeslabel,above] at (-.72,.18) {$R_1$};\node[xeslabel,above] at (.98,.18) {$R_2$};
    \node[xessmall] at (0,-.72) {串联：电流相同};
  \end{scope}
  \begin{scope}[shift={(3.55,0)}]
    \draw[circuit] (-2.05,0)--(-1.55,0) (-1.55,0)--(-1.55,1.0) (-1.55,0)--(-1.55,-1.0);
    \draw[circuit,decorate,decoration={zigzag,segment length=2.8mm,amplitude=1.05mm}] (-1.55,1.0)--(1.55,1.0);
    \draw[circuit,decorate,decoration={zigzag,segment length=2.8mm,amplitude=1.05mm}] (-1.55,-1.0)--(1.55,-1.0);
    \draw[circuit] (1.55,1.0)--(1.55,0) (1.55,-1.0)--(1.55,0) (1.55,0)--(2.05,0);
    \node[xeslabel,above] at (0,1.12) {$R_1$};\node[xeslabel,below] at (0,-1.12) {$R_2$};
    \node[xessmall] at (0,-1.68) {并联：电压相同};
  \end{scope}
\end{tikzpicture}
\end{center}}

\newcommand{\networkTermsFigure}{%
\begin{center}
\begin{tikzpicture}[x=1cm,y=1cm]
  \coordinate (A) at (0,2.0);\coordinate (B) at (-2.0,.62);\coordinate (C) at (-1.25,-1.72);
  \coordinate (D) at (1.25,-1.72);\coordinate (E) at (2.0,.62);\coordinate (O) at (0,0);
  \draw[circuit] (A)--(B)--(C)--(D)--(E)--cycle;
  \draw[circuit] (A)--(O)--(B) (O)--(C) (O)--(D) (O)--(E);
  \foreach \P in {A,B,C,D,E,O}{\node[nodepoint] at (\P) {};}
  \foreach \X/\Y in {A/B,B/C,C/D,D/E,E/A,A/O,O/B,O/C,O/D,O/E}{
    \path (\X)--(\Y) node[element,pos=.48,sloped] {};}
  \draw[Gold!80!black,densely dashed,line width=.7pt,rounded corners=7pt]
    (-.15,1.48)--(-.65,.58)--(0,0)--(1.05,.35)--(.85,1.20)--cycle;
  \node[xessmall,text=Gold!70!black] at (.60,.68) {回路};
  \draw[grayarrow] (-3.45,1.65)--(-2.05,.68);\node[xeslabel,anchor=east] at (-3.52,1.65) {节点};
  \draw[grayarrow] (-3.45,-.25)--($(B)!.60!(C)$);\node[xeslabel,anchor=east] at (-3.52,-.25) {支路};
  \draw[grayarrow] (3.55,1.35)--($(E)!.48!(A)$);\node[xeslabel,anchor=west] at (3.62,1.35) {元件};
  \node[xessmall] at (0,-2.32) {节点由导线连续区域确定；两个节点之间是一条支路};
\end{tikzpicture}
\end{center}}

\newcommand{\kclKvlFigure}{%
\begin{center}
\begin{tikzpicture}[x=.95cm,y=.95cm]
  \begin{scope}[shift={(-3.9,0)}]
    \coordinate (N1) at (0,1.85);\coordinate (N2) at (-1.70,-1.15);\coordinate (N4) at (1.70,-1.15);\coordinate (N3) at (0,-.05);
    \draw[circuit] (N1)--(N2)--(N4)--cycle (N1)--(N3)--(N2) (N3)--(N4);
    \foreach \P in {N1,N2,N3,N4}{\node[nodepoint] at (\P) {};}
    \foreach \X/\Y in {N1/N2,N1/N4,N2/N4,N1/N3,N3/N2,N3/N4}{\path (\X)--(\Y) node[element,pos=.48,sloped] {};}
    \draw[xesarrow] ($(N1)!.09!(N2)$)--($(N1)!.29!(N2)$);
    \node[xestiny,anchor=east] at (-.47,1.18) {$i_1$};
    \draw[xesarrow] ($(N1)!.08!(N3)$)--($(N1)!.28!(N3)$);
    \node[xestiny,anchor=west] at (.12,1.50) {$i_2$};
    \draw[xesarrow] ($(N1)!.09!(N4)$)--($(N1)!.29!(N4)$);
    \node[xestiny,anchor=west] at (.47,1.18) {$i_3$};
    \draw[xesarrow] ($(N3)!.08!(N2)$)--($(N3)!.29!(N2)$);
    \node[xestiny,anchor=east] at (-.30,-.22) {$i_4$};
    \draw[xesarrow] ($(N3)!.08!(N4)$)--($(N3)!.29!(N4)$);
    \node[xestiny,anchor=west] at (.30,-.22) {$i_5$};
    \draw[xesarrow] ($(N2)!.07!(N4)$)--($(N2)!.29!(N4)$);
    \node[xestiny,below] at (-1.05,-1.20) {$i_6$};
    \node[xeslabel,above] at (N1) {$N_1$};\node[xeslabel,below left] at (N2) {$N_2$};
    \node[xeslabel,above=2pt] at (N3) {$N_3$};\node[xeslabel,below right] at (N4) {$N_4$};
    \node[xessmall] at (0,-2.05) {KCL：先规定支路电流方向};
  \end{scope}
  \begin{scope}[shift={(3.9,0)}]
    \coordinate (N1) at (0,1.85);\coordinate (N2) at (-1.70,-1.15);\coordinate (N4) at (1.70,-1.15);\coordinate (N3) at (0,-.05);
    \draw[circuit] (N1)--(N2)--(N4)--cycle (N1)--(N3)--(N2) (N3)--(N4);
    \foreach \P in {N1,N2,N3,N4}{\node[nodepoint] at (\P) {};}
    \foreach \X/\Y in {N1/N2,N1/N4,N2/N4,N1/N3,N3/N2,N3/N4}{\path (\X)--(\Y) node[element,pos=.48,sloped] {};}
    \node[xestiny,anchor=east] at (-1.13,.63) {$v_1$};
    \node[xestiny,anchor=west] at (1.13,.63) {$v_3$};
    \node[xestiny,anchor=west] at (.18,1.03) {$v_2$};
    \node[xestiny,anchor=east] at (-.72,-.30) {$v_4$};
    \node[xestiny,anchor=west] at (.72,-.30) {$v_5$};
    \node[xestiny,below] at (0,-1.46) {$v_6$};
    \node[xestiny] at (-.39,.88) {$+$};\node[xestiny] at (-.94,-.08) {$-$};
    \node[xestiny] at (-.18,1.25) {$+$};\node[xestiny] at (-.18,.61) {$-$};
    \node[xestiny] at (.39,.88) {$+$};\node[xestiny] at (.94,-.08) {$-$};
    \node[xestiny] at (-.39,-.44) {$+$};\node[xestiny] at (-.94,-.80) {$-$};
    \node[xestiny] at (.39,-.44) {$+$};\node[xestiny] at (.94,-.80) {$-$};
    \node[xestiny] at (-.61,-.98) {$+$};\node[xestiny] at (.61,-.98) {$-$};
    \draw[Gold!80!black,-{Latex[length=1.5mm]},line width=.65pt]
      (-.66,.58) arc[start angle=120,end angle=430,radius=.20];
    \node[xestiny,text=Gold!70!black] at (-.58,.20) {$L_1$};
    \draw[Gold!80!black,-{Latex[length=1.5mm]},line width=.65pt]
      (.46,.58) arc[start angle=60,end angle=-250,radius=.20];
    \node[xestiny,text=Gold!70!black] at (.58,.20) {$L_2$};
    \draw[Gold!80!black,-{Latex[length=1.5mm]},line width=.65pt]
      (-.20,-.70) arc[start angle=170,end angle=-140,radius=.20];
    \node[xestiny,text=Gold!70!black] at (0,-.78) {$L_3$};
    \draw[Gold!80!black,-{Latex[length=1.7mm]},line width=.75pt]
      (0,2.28) arc[start angle=90,end angle=430,x radius=2.25,y radius=2.02];
    \node[xestiny,text=Gold!70!black] at (1.92,1.52) {$L_4$};
    \node[xessmall] at (0,-2.05) {KVL：沿所选回路累计支路电压};
  \end{scope}
\end{tikzpicture}
\end{center}}

\newcommand{\nodeMethodFigure}{%
\begin{center}
\begin{tikzpicture}[x=1cm,y=1cm]
  \coordinate (a) at (-2.4,1.35);\coordinate (b) at (1.25,1.35);\coordinate (c) at (1.25,-1.35);
  \draw[circuit] (a)--(-1.25,1.35);
  \draw[circuit,decorate,decoration={zigzag,segment length=3mm,amplitude=1.05mm}] (-1.25,1.35)--(.15,1.35);
  \draw[circuit] (.15,1.35)--(b);
  \draw[circuit] (b)--(1.25,.40);
  \draw[circuit,decorate,decoration={zigzag,segment length=3mm,amplitude=1.05mm}] (1.25,.40)--(1.25,-.78);
  \draw[circuit] (1.25,-.78)--(c)--(-2.4,-1.35);
  \draw[circuit] (-2.4,-1.35)--(-2.4,-.34);
  \draw[circuit] (-2.63,-.34)--(-2.17,-.34) (-2.78,.02)--(-2.02,.02);
  \draw[circuit] (-2.4,.02)--(a);
  \draw[circuit] (b)--(3.0,1.35)--(3.0,.34);
  \draw[circuit] (3.0,-.34)--(3.0,-1.35)--(c);
  \draw[circuit] (3.0,0) circle (.34);
  \draw[xesarrow] (3.0,-.20)--(3.0,.20);
  \foreach \P in {a,b,c}{\node[nodepoint] at (\P) {};}
  \node[xeslabel,above] at (a) {$a$};\node[xeslabel,above] at (b) {$b$};\node[xeslabel,below] at (c) {$c$};
  \node[xeslabel,above] at (-.55,1.48) {$R_1$};\node[xeslabel,right] at (1.42,-.18) {$R_2$};
  \node[xeslabel,left] at (-2.82,-.15) {$E_0$};\node[xeslabel,right] at (3.40,0) {$I_0$};
  \draw[circuit] (c)--(1.25,-1.65) (.92,-1.65)--(1.58,-1.65) (1.02,-1.82)--(1.48,-1.82) (1.12,-1.98)--(1.38,-1.98);
  \node[xessmall] at (0,-2.42) {取 $c$ 为参考地后，只需求节点 $b$ 的电势};
\end{tikzpicture}
\end{center}}

\newcommand{\tolmanRingFigure}{%
\begin{center}
\begin{tikzpicture}[x=1cm,y=1cm]
  \draw[xesblue,line width=1.2pt] (0,0) ellipse (2.55 and 1.12);
  \draw[xesblue,line width=.75pt] (0,0) ellipse (2.30 and .92);
  \draw[xesline] (0,-1.55)--(0,1.55);
  \draw[xesarrow] (-.55,1.45) arc[start angle=160,end angle=25,x radius=.65,y radius=.30];
  \node[xeslabel,above] at (0,1.52) {$\omega$};
  \draw[xesarrow] (2.08,.52)..controls(1.55,.85)and(.95,1.02)..(.35,1.08);
  \node[xeslabel,above] at (1.25,.98) {$I(t)$};
  \draw[RuleGray,dashed] (0,0)--(2.40,.37);\node[xeslabel,above] at (1.42,.27) {$r$};
  \node[xessmall] at (0,-1.98) {晶格骤停后，电子因惯性暂时保持切向速度};
\end{tikzpicture}
\end{center}}

\newcommand{\transientCurrentFigure}{%
\begin{center}
\begin{tikzpicture}[x=1cm,y=1cm]
  \draw[xesline,-{Latex[length=1.8mm]}] (0,0)--(5.2,0) node[xeslabel,right] {$t$};
  \draw[xesline,-{Latex[length=1.8mm]}] (0,0)--(0,2.55) node[xeslabel,above] {$I$};
  \draw[RuleGray,densely dashed] (0,2.0)--(4.75,2.0);
  \node[xessmall,anchor=east] at (-.12,2.0) {$U/R$};
  \draw[xesblue] plot[domain=0:4.7,samples=80] (\x,{2*(1-exp(-\x/1.0))});
  \draw[<->,Gold!75!black,line width=.65pt] (0,-.42)--(1.0,-.42);
  \node[xestiny,text=Gold!70!black,below] at (.5,-.46) {$\tau$};
\end{tikzpicture}
\end{center}}

\newcommand{\rotatingChargedRingFigure}{%
\begin{center}
\begin{tikzpicture}[x=1cm,y=1cm]
  \draw[xesblue,line width=1.2pt] (0,0) ellipse (2.55 and 1.03);
  \draw[xesblue,line width=.72pt] (0,0) ellipse (2.30 and .85);
  \draw[xesline] (0,-1.50)--(0,1.50);
  \draw[xesarrow] (-.55,1.40) arc[start angle=160,end angle=25,x radius=.65,y radius=.30];
  \node[xeslabel,above] at (0,1.48) {$\omega$};
  \draw[xesarrow] (.48,1.02)..controls(1.1,.96)and(1.65,.83)..(2.15,.50);
  \node[xeslabel,above] at (1.37,.95) {$I$};
  \draw[RuleGray,dashed] (0,0)--(2.38,.35);\node[xeslabel,above] at (1.42,.24) {$R$};
  \draw[grayarrow] (-3.05,.62)--(-2.34,.51);\node[xeslabel,anchor=east] at (-3.12,.62) {$Q$};
  \node[xeslabel,below left] at (0,0) {$O$};
\end{tikzpicture}
\end{center}}

\newcommand{\vacuumDiodeFigure}{%
\begin{center}
\begin{tikzpicture}[x=1cm,y=1cm]
  \draw[xesblue,line width=1.25pt] (-2.35,-.75)--(-2.35,.75) (2.35,-.75)--(2.35,.75);
  \draw[<->,Gold!75!black,line width=.65pt] (-2.30,1.52)--(2.30,1.52);
  \node[xestiny,text=Gold!70!black] at (0,1.76) {$L$};
  \node[xeslabel,above] at (-2.35,.82) {阴极 $C$};\node[xeslabel,above] at (2.35,.82) {阳极 $A$};
  \draw[xesarrow] (-1.75,0)--(1.75,0);\node[xessmall,above] at (0,.05) {电子运动};
  \draw[grayarrow] (2.95,.45)--(2.40,.45);\node[xeslabel,anchor=west] at (3.02,.45) {$I$};
  \draw[xesline,-{Latex[length=1.6mm]}] (-2.60,-1.05)--(2.80,-1.05) node[xeslabel,right] {$x$};
  \node[xestiny,below] at (-2.35,-1.06) {$0$};\node[xestiny,below] at (2.35,-1.06) {$L$};
\end{tikzpicture}
\end{center}}

\newcommand{\childLangmuirFigure}{%
\begin{center}
\begin{tikzpicture}[x=1cm,y=1cm]
  \draw[xesline,-{Latex[length=1.7mm]}] (0,0)--(5.15,0) node[xeslabel,right] {$x$};
  \draw[xesline,-{Latex[length=1.7mm]}] (0,0)--(0,2.65) node[xeslabel,above] {$\varphi$};
  \draw[xesblue] plot[domain=0:4.7,samples=80] (\x,{2.25*(\x/4.7)^(4/3)});
  \draw[RuleGray,densely dashed] (4.7,0)--(4.7,2.25)--(0,2.25);
  \node[xestiny,below] at (4.7,0) {$L$};\node[xestiny,anchor=east] at (-.10,2.25) {$U$};
  \node[xessmall,anchor=west] at (1.00,1.82) {$\varphi(x)=U(x/L)^{4/3}$};
\end{tikzpicture}
\end{center}}

\newcommand{\conductingShellFigure}{%
\begin{center}
\begin{tikzpicture}[x=1cm,y=1cm]
  \fill[DeepBlue!8] (0,0) circle (1.85);\fill[white] (0,0) circle (.82);
  \draw[xesblue] (0,0) circle (1.85) (0,0) circle (.82);
  \draw[xesline] (0,0)--(.58,-.58);\node[xeslabel,below right] at (.62,-.62) {$a$};
  \draw[xesline] (0,0)--(0,1.85);\node[xeslabel,left] at (0,1.25) {$b$};
  \draw[xesarrow] (1.05,.25)--(1.62,.39);\node[xessmall,anchor=west] at (1.72,.44) {$\boldsymbol j\parallel\boldsymbol E$};
  \node[xessmall] at (0,-2.30) {$\sigma(r)=K E(r)$，电流在任意同心球面上相同};
\end{tikzpicture}
\end{center}}

\newcommand{\ionChamberFigure}{%
\begin{center}
\begin{tikzpicture}[x=1cm,y=1cm]
  \draw[xesblue,line width=1.2pt] (-3.5,1.45)--(3.5,1.45) (-3.5,-1.45)--(3.5,-1.45);
  \draw[xesline] (-3.5,1.45)--(-3.5,-1.45) (3.5,1.45)--(3.5,-1.45);
  \node[xeslabel,above] at (-1.90,1.48) {阴极};\node[xeslabel,below] at (-1.90,-1.48) {阳极};
  \draw[circuit] (1.90,1.45)--(1.90,2.05)--(3.15,2.05);\node[xeslabel,right] at (3.15,2.05) {$-U_0$};
  \draw[circuit] (1.90,-1.45)--(1.90,-2.05);
  \draw[circuit,decorate,decoration={zigzag,segment length=2.7mm,amplitude=1.0mm}] (1.90,-2.05)--(1.90,-2.95);
  \draw[circuit] (1.90,-2.95)--(1.90,-3.20);
  \draw[circuit] (1.55,-3.20)--(2.25,-3.20) (1.66,-3.37)--(2.14,-3.37) (1.77,-3.54)--(2.03,-3.54);
  \node[xeslabel,right] at (2.10,-2.50) {$R$};\draw[circuit] (1.90,-1.70)--(3.20,-1.70);
  \node[xeslabel,right] at (3.20,-1.70) {信号};
  \draw[xesline,-{Latex[length=1.7mm]}] (-4.65,-1.45)--(-4.65,2.05) node[xeslabel,above] {$x$};
  \draw[<->,Gold!75!black,line width=.65pt] (-3.78,-1.42)--(-3.78,1.42);
  \node[xestiny,text=Gold!70!black,anchor=east] at (-3.92,0) {$D$};
  \draw[<->,Gold!75!black,line width=.65pt] (-3.25,-1.40)--(-3.25,-.15);
  \node[xestiny,text=Gold!70!black,anchor=east] at (-3.36,-.80) {$x_0$};
  \draw[xesarrow] (-2.85,-.14)--(2.85,-.14);\node[xessmall,above] at (1.5,.75) {入射粒子径迹};
  \foreach \x in {-2.3,-1.65,-.95,-.25,.45,1.15,1.85,2.50}{
    \node[xestiny,text=Gold!70!black] at (\x,.02) {$+$};\node[xestiny,text=DeepBlue] at (\x,-.29) {$-$};}
  \draw[grayarrow] (-.42,.08)--(-.42,.78);\node[xestiny,anchor=west] at (-.34,.63) {离子 $v^+$};
  \draw[grayarrow] (.42,-.32)--(.42,-.98);\node[xestiny,anchor=west] at (.50,-.88) {电子 $v^-$};
  \node[xessmall] at (0,2.58) {气密容器内充工作气体；电场把电子与正离子分别引向两极};
\end{tikzpicture}
\end{center}}

\newcommand{\ionSignalFigure}{%
\begin{center}
\begin{tikzpicture}[x=1cm,y=1cm]
  \draw[xesline,-{Latex[length=1.7mm]}] (0,0)--(6.1,0) node[xeslabel,right] {$t$};
  \draw[xesline,-{Latex[length=1.7mm]}] (0,0)--(0,2.65) node[xeslabel,above] {$|\Delta U|$};
  \draw[xesblue] (0,0)--(1.25,1.25)--(4.70,2.25);
  \draw[RuleGray,densely dashed] (1.25,0)--(1.25,1.25) (4.70,0)--(4.70,2.25);
  \node[xestiny,below] at (1.25,0) {$t_e=x_0/v^-$};
  \node[xestiny,below] at (4.70,0) {$t_i=(D-x_0)/v^+$};
  \node[xessmall] at (3.10,2.55) {电子收集完后，离子继续形成较慢斜坡};
\end{tikzpicture}
\end{center}}

\newcommand{\ringResistorFigure}{%
\begin{center}
\begin{tikzpicture}[x=1cm,y=1cm]
  \draw[xesblue,line width=1.1pt] (0,0) circle (1.55);
  \coordinate (a) at (-1.55,0);\coordinate (b) at ({1.55*cos(-35)},{1.55*sin(-35)});
  \node[nodepoint] at (a) {};\node[nodepoint] at (b) {};
  \draw[circuit] (a)--(-2.65,0);\draw[circuit] (b)--($(b)+(1.0,-.70)$);
  \draw[RuleGray,dashed] (0,0)--(a) (0,0)--(b);
  \node[xeslabel,above left] at (a) {$a$};\node[xeslabel,below right] at (b) {$b$};
  \node[xeslabel,below left] at (0,0) {$O$};\node[xeslabel] at (-.45,.90) {$R_1$};\node[xeslabel] at (.72,-.80) {$R_2$};
  \node[xessmall] at (0,-2.10) {$R_1+R_2=R$，两段圆弧在端点 $a,b$ 之间并联};
\end{tikzpicture}
\end{center}}

\newcommand{\dualSourceFigure}{%
\begin{center}
\begin{tikzpicture}[x=1cm,y=1cm]
  \coordinate (L) at (-2.65,0);\coordinate (R) at (2.65,0);
  \draw[circuit] (L)--(-1.50,0);
  \draw[circuit,decorate,decoration={zigzag,segment length=3mm,amplitude=1.05mm}] (-1.50,0)--(1.20,0);
  \draw[circuit] (1.20,0)--(R);
  \node[xeslabel,above] at (0,0.25) {$R$};
  \draw[circuit] (L)--(-2.65,-1.20)--(-.55,-1.20);
  \draw[circuit] (-.55,-1.60)--(-.55,-.80) (-.15,-1.48)--(-.15,-.92);
  \draw[circuit] (-.15,-1.20)--(.35,-1.20);
  \draw[circuit,decorate,decoration={zigzag,segment length=2.8mm,amplitude=.9mm}] (.35,-1.20)--(1.70,-1.20);
  \draw[circuit] (1.70,-1.20)--(2.65,-1.20)--(R);
  \node[xeslabel,below] at (-.35,-1.62) {$\varepsilon_1$};\node[xeslabel,above] at (1.02,-1.03) {$r_1$};
  \draw[circuit] (L)--(-2.65,-2.40)--(-.55,-2.40);
  \draw[circuit] (-.55,-2.80)--(-.55,-2.00) (-.15,-2.68)--(-.15,-2.12);
  \draw[circuit] (-.15,-2.40)--(.35,-2.40);
  \draw[circuit,decorate,decoration={zigzag,segment length=2.8mm,amplitude=.9mm}] (.35,-2.40)--(1.70,-2.40);
  \draw[circuit] (1.70,-2.40)--(2.65,-2.40)--(R);
  \node[xeslabel,below] at (-.35,-2.82) {$\varepsilon_2$};\node[xeslabel,above] at (1.02,-2.23) {$r_2$};
  \node[nodepoint] at (L) {};\node[nodepoint] at (R) {};
  \draw[xesarrow] (-1.9,.45)--(-.85,.45);\node[xeslabel,above] at (-1.38,.48) {$I$};
  \node[xessmall] at (0,-3.30) {两个实际电源同极性并联后向负载 $R$ 供电};
\end{tikzpicture}
\end{center}}

% Source question conditions restored, analysis retained separately.
\newcommand{\emRotatingRingQuestion}{%
\begin{center}
\begin{tikzpicture}[x=1cm,y=1cm]
  \draw[xesblue,line width=1.2pt] (0,0) ellipse (2.55 and 1.03);
  \draw[xesblue,line width=.72pt] (0,0) ellipse (2.30 and .85);
  \draw[xesline] (0,-1.50)--(0,1.50);
  \draw[xesarrow] (-.55,1.40) arc[start angle=160,end angle=25,x radius=.65,y radius=.30];
  \node[xeslabel,above] at (0,1.48) {$\omega$};
  \draw[RuleGray,dashed] (0,0)--(2.38,.35);\node[xeslabel,above] at (1.42,.24) {$R$};
  \draw[grayarrow] (-3.05,.62)--(-2.34,.51);\node[xeslabel,anchor=east] at (-3.12,.62) {$Q$};
  \node[xeslabel,below left] at (0,0) {$O$};
\end{tikzpicture}
\end{center}}

% Source question conditions restored, analysis retained separately.
\newcommand{\emConductingShellQuestion}{%
\begin{center}
\begin{tikzpicture}[x=1cm,y=1cm]
  \fill[DeepBlue!8] (0,0) circle (1.85);\fill[white] (0,0) circle (.82);
  \draw[xesblue] (0,0) circle (1.85) (0,0) circle (.82);
  \draw[xesline] (0,0)--(.58,-.58);\node[xeslabel,below right] at (.62,-.62) {$a$};
  \draw[xesline] (0,0)--(0,1.85);\node[xeslabel,left] at (0,1.25) {$b$};
\end{tikzpicture}
\end{center}}

% Source question conditions restored, analysis retained separately.
\newcommand{\emIonChamberQuestion}{%
\begin{center}
\begin{tikzpicture}[x=1cm,y=1cm]
  \draw[xesblue,line width=1.2pt] (-3.5,1.45)--(3.5,1.45) (-3.5,-1.45)--(3.5,-1.45);
  \draw[xesline] (-3.5,1.45)--(-3.5,-1.45) (3.5,1.45)--(3.5,-1.45);
  \node[xeslabel,above] at (-1.90,1.48) {阴极};\node[xeslabel,below] at (-1.90,-1.48) {阳极};
  \draw[circuit] (1.90,1.45)--(1.90,2.05)--(3.15,2.05);\node[xeslabel,right] at (3.15,2.05) {$-U_0$};
  \draw[circuit] (1.90,-1.45)--(1.90,-2.05);
  \draw[circuit,decorate,decoration={zigzag,segment length=2.7mm,amplitude=1.0mm}] (1.90,-2.05)--(1.90,-2.95);
  \draw[circuit] (1.90,-2.95)--(1.90,-3.20);
  \draw[circuit] (1.55,-3.20)--(2.25,-3.20) (1.66,-3.37)--(2.14,-3.37) (1.77,-3.54)--(2.03,-3.54);
  \node[xeslabel,right] at (2.10,-2.50) {$R$};\draw[circuit] (1.90,-1.70)--(3.20,-1.70);
  \node[xeslabel,right] at (3.20,-1.70) {信号};
  \draw[xesline,-{Latex[length=1.7mm]}] (-4.65,-1.45)--(-4.65,2.05) node[xeslabel,above] {$x$};
  \draw[<->,Gold!75!black,line width=.65pt] (-3.78,-1.42)--(-3.78,1.42);
  \node[xestiny,text=Gold!70!black,anchor=east] at (-3.92,0) {$D$};
  \draw[<->,Gold!75!black,line width=.65pt] (-3.25,-1.40)--(-3.25,-.15);
  \node[xestiny,text=Gold!70!black,anchor=east] at (-3.36,-.80) {$x_0$};
  \draw[xesarrow] (-2.85,-.14)--(2.85,-.14);\node[xessmall,above] at (1.5,.75) {入射粒子径迹};
  \foreach \x in {-2.3,-1.65,-.95,-.25,.45,1.15,1.85,2.50}{
    \node[xestiny,text=Gold!70!black] at (\x,.02) {$+$};\node[xestiny,text=DeepBlue] at (\x,-.29) {$-$};}
\end{tikzpicture}
\end{center}}

% Source question conditions restored, analysis retained separately.
\newcommand{\emRingResistorQuestion}{%
\begin{center}
\begin{tikzpicture}[x=1cm,y=1cm]
  \draw[xesblue,line width=1.1pt] (0,0) circle (1.55);
  \coordinate (a) at (-1.55,0);\coordinate (b) at ({1.55*cos(-35)},{1.55*sin(-35)});
  \node[nodepoint] at (a) {};\node[nodepoint] at (b) {};
  \draw[circuit] (a)--(-2.65,0);\draw[circuit] (b)--($(b)+(1.0,-.70)$);
  \draw[RuleGray,dashed] (0,0)--(a) (0,0)--(b);
  \node[xeslabel,above left] at (a) {$a$};\node[xeslabel,below right] at (b) {$b$};
  \node[xeslabel,below left] at (0,0) {$O$};\node[xeslabel] at (-.45,.90) {$R_1$};\node[xeslabel] at (.72,-.80) {$R_2$};
\end{tikzpicture}
\end{center}}


在最一般的意义下，电流是大量带电粒子沿确定方向发生整体定向运动的宏观表现。电能的输运和使用
已经成为现代社会的基础；如果所有电力驱动的设备突然停止工作，生活与生产会立即受到全面影响。
本章从电流的定义出发，讨论物质导电的经典图像、元件的伏安特性、欧姆定律（Ohm's law）以及一般直流网络的
基尔霍夫定律（Kirchhoff's circuit laws）。本章后续内容还将在这些基础上研究更复杂的电阻网络。

导体和绝缘体的关键差别，在于物质内部是否存在能够在外场作用下长距离迁移的带电粒子，即载流子。
金属中的载流子主要是电子；电解液中正、负离子都可参与导电；半导体中则可能同时存在电子和空穴。
载流子的浓度、质量、电荷量和碰撞性质，以及温度、外电场、光照等外部条件，共同决定材料的导电性。
必须强调：金属导电的精确理论属于量子物理。下面的经典模型不能从第一性原理给出全部材料参数，
但能把宏观规律与载流子的运动联系起来，并说明经典图像在哪些条件下有效、又会在何处失效。


\subsection{电流强度和电流密度}


电流并不只存在于接有电源的闭合回路中；只要带电粒子发生有方向的输运，就形成电流。取一个有向截面，
若在时间间隔 $\Delta t$ 内按规定正方向通过它的净电量为 $\Delta Q$，平均电流强度定义为
\begin{equation}
  \boxed{I=\frac{\Delta Q}{\Delta t}}.
\end{equation}
当时间间隔趋于零，瞬时电流是通过电量 $Q(t)$ 对时间的导数：
\begin{equation}
  \boxed{I(t)=\lim_{\Delta t\to0}\frac{\Delta Q}{\Delta t}
  =\frac{\mathrm dQ}{\mathrm dt}}.
\end{equation}
一般情况下 $I$ 随时间变化；若它不随时间变化，称为稳恒电流。反过来，已知电流随时间的关系后，
从 $t_1$ 到 $t_2$ 通过截面的净电量为
\begin{equation}
  \boxed{Q(t_1\to t_2)=\int_{t_1}^{t_2}I(t)\,\mathrm dt}.
\end{equation}
对于恒定电流，上式退化为 $Q=I(t_2-t_1)$。电流的国际单位是安培，
$1\,\mathrm A=1\,\mathrm C\,\mathrm s^{-1}$。

设导线横截面积为 $A$，单位体积内有 $n$ 个同种载流子，每个载流子的电量为 $q$，沿导线的平均定向
漂移速度为 $v_d$。在 $\Delta t$ 内能够穿过截面的载流子原来位于长为 $v_d\Delta t$ 的柱体中，
其数目为 $nAv_d\Delta t$，因而通过的电量是 $nqAv_d\Delta t$。代入定义得到
\begin{equation}
  \boxed{I=nqAv_d}.
\end{equation}
式中 $q$ 是带符号的电荷量，$v_d$ 是沿所选截面法向的速度分量。对电子，电荷与沿常规电流方向的
速度分量均为负，乘积为正。若只讨论大小，则写 $|I|=n|q|A|v_d|$，不可只对电荷取绝对值。

真实导线中的漂移速度通常很小。以横截面积 $A=1\,\mathrm{mm^2}$ 的纯银导线为例，若每个银原子提供
一个导电电子，则 $n\simeq6\times10^{28}\,\mathrm{m^{-3}}$。通有 $1\,\mathrm A$ 电流时
\begin{equation}
  \boxed{v_d=\frac{I}{neA}\approx1.0\times10^{-4}\,\mathrm{m\,s^{-1}}}.
\end{equation}
这个速度比金属中电子本身的无规则运动速度小许多个数量级。因此“电子像小球一样沿导线迅速跑到用电器”
并不是正确图像：电路中建立电场的过程很快，但单个电子的净漂移极慢。
金属电子的无规则运动速率不能用经典能量均分定理准确求得，必须由量子统计决定；其典型高速电子可达到
约光速的百分之一，而上面算出的净漂移仍只有 $10^{-4}\,\mathrm{m\,s^{-1}}$ 量级。

若截面较大，不同位置的载流子漂移可能并不相同，单一的 $I$ 不能描述局部流动。定义电流密度矢量
\begin{equation}
  \boxed{\boldsymbol j=nq\boldsymbol v_d}.
\end{equation}
它的方向按正电荷运动方向规定。对于电子，$q=-e$，所以 $\boldsymbol j$ 与电子漂移速度方向相反；
若有多种载流子同时运动，总电流密度是各类载流子贡献之和。穿过任意有向曲面 $S$ 的电流为
\[
  I=\iint_S\boldsymbol j\cdot\mathrm d\boldsymbol S.
\]
这正像流体的体积流量等于速度场在截面上的通量，只不过电流计算的是电荷通量。
例如河水中央常比近岸处流得快，必须先描述各点流速才能求总流量；非均匀导体中的电流也完全类似。


\subsection{导电的微观机制——简单模型}


先构造一个纯粹唯象的“黏滞阻力”模型。设载流子定向运动时受到与漂移速度成正比、方向相反的平均阻力
$\boldsymbol F_r=-\alpha\boldsymbol v_d$。恒定电场中达到稳定漂移后，电场力与阻力平衡：
\[
  q\boldsymbol E-\alpha\boldsymbol v_d=0,
  \qquad \boldsymbol v_d=\frac q\alpha\boldsymbol E.
\]
于是
\[
  \boldsymbol j=nq\boldsymbol v_d=\frac{nq^2}{\alpha}\boldsymbol E.
\]
与微观欧姆定律比较可得
\begin{equation}
  \boxed{\sigma=\frac{nq^2}{\alpha}}.
\end{equation}
这个模型把极复杂的散射过程压缩成一个平均阻力系数。它虽然不是真实金属中电子运动的精确图像，却能
说明为什么稳定状态下 $\boldsymbol j$ 与 $\boldsymbol E$ 同向且成正比。

它还提示了暂态过程：电压突然改变时，载流子具有惯性，漂移速度不能无延迟地跳变。若保留质量 $m$，
运动方程为 $m\dot{\boldsymbol v}_d=q\boldsymbol E-\alpha\boldsymbol v_d$，电流会在特征时间
$\tau=m/\alpha$ 内趋近新的稳态值。通常金属中的 $\tau$ 极短，所以普通低频电路把这一过程近似为
瞬时完成；例~\ref{ch21:ex:s36-5-2} 将把它定量写出。


% EM teaching enrichment 20260927 supplementary

\paragraph{从建立过程看稳定电流}
这个模型也给出电场接通后的响应。初始漂移速度为零时，
\[
 m\frac{\dd\vect v_d}{\dd t}=q\vect E-\alpha\vect v_d,
 \qquad \vect v_d(t)=\frac{q\vect E}{\alpha}(1-\mathrm e^{-t/\tau}),\quad \tau=\frac m\alpha.
\]
乘以 $nq$ 后得到 $\vect j(t)=\sigma\vect E(1-\mathrm e^{-t/\tau})$，其中 $\sigma=nq^2/\alpha$。
达到 $t=\tau$ 时电流约为最终值的 $63\%$，约数个 $\tau$ 后才可用稳态欧姆定律。
负载流子的漂移方向与电场相反，但 $q\vect v_d$ 与电场同向，故电导率仍为正。
这一时间尺度描述载流子的局部松弛，不能把它当作信号沿整根导线传播的时间。


\subsection{电子的扩散和漂移}


金属中的电子并非只做缓慢定向运动。它们本来就在晶体中高速无规则运动，外电场只在这种运动上叠加
一个很小的平均漂移。电子在两次散射之间近似自由运动，受电场加速；散射会把先前积累的定向速度随机化，
大量电子的统计平均才形成稳定电流。

\par\noindent\begin{minipage}{\linewidth}
\drudeFigure
\diagramcaption{导体内电子的无规则运动与定向漂移。蓝色曲线表示一条示意轨迹，金色圆点表示晶格散射中心。}\label{ch21:fig:s36-5-6}
\end{minipage}\par\medskip

设自由电子数密度为 $n$、电子电荷量的绝对值为 $e$、质量为 $m$，两次有效散射之间的平均时间为
$\tau$。在弱电场中，电子由电场获得的平均漂移速度大小约为 $v_d=eE\tau/m$。因此
\begin{equation}
  \boxed{\boldsymbol j=\sigma\boldsymbol E,\qquad
  \sigma=\frac{ne^2\tau}{m}}.
\end{equation}
这就是经典德鲁德模型给出的电导率。若从一次碰撞后零定向速度开始计算到下一次碰撞，简单的匀加速
平均会出现 $1/2$；但真实碰撞时间服从分布，玻尔兹曼方程的弛豫时间近似给出上式的标准系数。

该模型也帮助理解金属电阻率随温度升高而增大的现象。金属中载流子浓度 $n$ 大致不变；升温增强晶格
振动，使散射更频繁、$\tau$ 变小，于是 $\sigma$ 减小而 $\rho$ 增大。半导体则不同：升温会显著增加
载流子浓度，$n$ 增大的效果常超过 $\tau$ 减小的效果，因此其电阻通常随温度升高而减小。
在较粗略的经典估计中，也可暂把平均自由程 $\ell$ 看作不变：温度升高使电子典型速率增大，
$\tau\sim\ell/\bar v$ 因而减小，仍得到同一趋势；但金属中更准确的解释需要量子统计与晶格散射理论。

推导上述电导率公式时隐含了几项条件：电子在一次自由飞行中由外场获得的定向速度远小于其无规则运动速度；
散射机制和平均碰撞时间不被外场显著改变；散射足够频繁；载流子浓度不随外场改变。任何一项遭到破坏，
伏安关系都可能偏离直线。稀薄气体在强电场中的放电就是典型例子：电子在两次碰撞之间获得很大能量，
可激发或电离气体分子，使载流子浓度发生剧烈变化，依次出现暗放电、辉光放电和电弧区。

\par\noindent\begin{minipage}{\linewidth}
\gasDischargeFigure
\diagramcaption{气体放电的典型电压--电流曲线。横轴为对数电流标度，各区段明显不满足线性欧姆定律。}\label{ch21:fig:s36-5-7}
\end{minipage}\par\medskip

经典模型仍不能精确描述金属电子，但能把实测电导率换算成有意义的微观尺度。以银为例，
$\sigma\simeq6.3\times10^7\,\mathrm{S\,m^{-1}}$，取
$n\simeq6.0\times10^{28}\,\mathrm{m^{-3}}$，由 $\sigma=ne^2\tau/m$ 得到
\begin{equation}
  \boxed{\tau=\frac{m\sigma}{ne^2}\approx3.8\times10^{-14}\,\mathrm s}.
\end{equation}
若电子典型无规则运动速率取 $\bar v\sim10^6\,\mathrm{m\,s^{-1}}$，平均自由程约为
\begin{equation}
  \boxed{\ell\sim\bar v\tau\approx3.8\times10^{-8}\,\mathrm m}.
\end{equation}
它远大于晶格中原子间的典型距离，说明电子并非每经过一个原子就发生一次经典小球式碰撞；这一事实正是
后来建立能带理论和量子输运理论的重要线索。例如取原子间距约 $0.3\,\mathrm{nm}$，
上述 $38\,\mathrm{nm}$ 自由程约为其 $127$ 倍。这与朴素的“逐原子碰撞”图像相差很大。


% EM teaching enrichment 20260927
\paragraph{浓度不均匀时还有扩散流}
无规则热运动在均匀密度下彼此抵消；若载流子数密度 $n$ 随位置变化，粒子通量可近似写成
$\boldsymbol\Gamma_{\rm diff}=-D\nabla n$，相应电流为 $q\boldsymbol\Gamma_{\rm diff}$。加上漂移项有
\[
 \boldsymbol j=nq\boldsymbol v_d-qD\nabla n.
\]
这里 $D$ 是扩散系数，单位 $\mathrm{m^2/s}$。对电子 $q=-e$，扩散电流方向与电子扩散方向相反。
稳态零电流可以来自漂移与扩散相互抵消，而非各项分别为零。因此含浓度梯度的区域不能只据 $\boldsymbol j=0$ 就断言 $\boldsymbol E=0$；前面的纯欧姆关系省略了这一驱动力。


% EM structure and teaching revision 20260928
\subsection{电流连续性与载流子建立时间}
宏观电流密度由各类载流子的定向运动叠加：$\boldsymbol j=\sum_s n_sq_s\boldsymbol v_s$。电子的平均漂移方向与常规电流方向相反，但其随机热运动通常远快于漂移；随机速度在平衡状态下的平均值为零。
取简单碰撞模型，电子电荷为 $-e$、平均弛豫时间为 $\tau$，恒定外场接通后
\[
 m\frac{\dd\boldsymbol v_d}{\dd t}=-e\boldsymbol E-\frac m\tau\boldsymbol v_d,
 \qquad\boldsymbol v_d(t)=-\frac{e\tau}{m}\boldsymbol E(1-e^{-t/\tau}).
\]
长时间后得到 $\boldsymbol j=ne^2\tau\boldsymbol E/m$。这说明欧姆关系是该线性模型的稳定响应，而非任何材料、任何时间下都成立的定义。
局部电荷守恒写为 $\partial\rho/\partial t+\nabla\cdot\boldsymbol j=0$。稳定电流满足 $\nabla\cdot\boldsymbol j=0$，但不要求电流密度处处相同；在变截面导体中 $j=I/S$，较细处的场和单位体积发热都可能更大。
实际建立过程还受电磁扰动传播和电路储能影响，因此微观碰撞时间不能直接替代整条电路的RC或LR时间常数。


\subsection{专题例题与方法应用}

\begin{origexamplebox}{【惯性维持的电流】}\label{ch21:ex:s36-5-1}
1917 年，托尔曼和斯图尔特的实验直接证明了电子在金属中参与导电。他们让一个铜环以很高的恒定
角速度绕主轴转动，此时环中没有相对晶格的电流；再让铜环突然停止，发现环内出现短暂电流。求从铜环
停止开始流过任一横截面的总电量。结果用环半径 $r$、初始角速度 $\omega$、电子电量绝对值 $e$、
电子质量 $m$ 和铜的电导率 $\sigma$ 表示。记导线横截面积为 $S$。
\end{origexamplebox}\par\nopagebreak[4]

\par\noindent\begin{minipage}{\linewidth}
\tolmanRingFigure
\diagramcaption{解图：铜环晶格骤停后，电子因惯性保持原切向速度；常规电流方向与电子运动方向相反。}
\end{minipage}\par\medskip

\begin{xesolution}
停止前，自由电子随晶格共同转动，其切向速度为 $r\omega$；相对于晶格没有漂移，所以没有电流。晶格
突然停止时，电子速度不能突变，刚停止后的相对漂移速度仍为 $r\omega$。用德鲁德弛豫模型描述碰撞：
\[
  m\frac{\mathrm dv}{\mathrm dt}=-\frac m\tau v,
  \qquad v(0)=r\omega,
\]
故
\[
  v(t)=r\omega e^{-t/\tau}.
\]
电子数密度为 $n$，于是常规电流大小
\[
  |I(t)|=neS\,r\omega e^{-t/\tau}.
\]
从 $t=0$ 积分到电流完全衰减：
\[
  |Q|=\int_0^\infty |I(t)|\,\mathrm dt=neSr\omega\tau.
\]
再利用德鲁德关系 $\sigma=ne^2\tau/m$ 消去 $n$ 与 $\tau$，得到
\[
  \boxed{|Q|=\frac{m\sigma S r\omega}{e}}.
\]
若把转动方向规定为正，由于电子带负电，常规电流以及带符号的 $Q$ 均沿反方向。
总电量与横截面积 $S$ 成正比；单位横截面积通过的电量为
$|Q|/S=m\sigma r\omega/e$。
\end{xesolution}



\begin{origexamplebox}{【暂态电流】}\label{ch21:ex:s36-5-2}
在一个阻值为 $R$ 的圆柱形电阻两端突然加上恒定电压 $U$。考虑载流子的惯性，求通过电阻的电流
随时间的变化关系 $I(t)$。
\end{origexamplebox}\par\nopagebreak[4]

\begin{xesolution}
设电子漂移速度的大小为 $v_d$，取常规电流方向为正。德鲁德模型的平均运动方程可写成
\[
  m\frac{\mathrm dv_d}{\mathrm dt}+\frac m\tau v_d=eE.
\]
圆柱导体内 $E=U/l$，而 $I=neSv_d$。把方程乘以 $neS$，并利用稳态时
$U/R=neS(eE\tau/m)$，可得电流满足
\[
  \tau\frac{\mathrm dI}{\mathrm dt}+I=\frac UR.
\]
突然加压前 $I(0)=0$，解这个一阶方程：
\[
  \boxed{I(t)=\frac UR\left(1-e^{-t/\tau}\right),\qquad t\ge0}.
\]

\par\noindent\begin{minipage}{\linewidth}
\transientCurrentFigure
\diagramcaption{解图：电流以时间常数 $\tau$ 从零趋近稳态值 $U/R$；普通电路中 $\tau$ 极短。}
\end{minipage}\par\medskip

欧姆定律 $I=U/R$ 是 $t\gg\tau$ 后的稳态结果。若理想化地令 $\tau\to0$，才会得到“加压瞬间电流
立刻跳到 $U/R$”的集中参数模型。
\end{xesolution}



\begin{origexamplebox}{【转动的带电环】}\label{ch21:ex:s36-5-3}
一个半径为 $R$、总电荷为 $Q$ 的刚性均匀带电圆环，绕其对称轴以角速度 $\omega$ 匀速转动。
这一过程可以等效为圆形导线中的稳恒电流，求等效电流 $I$。
\par\noindent\begin{minipage}{\linewidth}
\emRotatingRingQuestion
\figcap{刚性带电圆环及给定转动方向，等效电流待求。}\label{ch21:fig:em-source-thirty-sixth-section-3}
\end{minipage}\par\medskip
\end{origexamplebox}\par\nopagebreak[4]

\par\noindent\begin{minipage}{\linewidth}
\rotatingChargedRingFigure
\diagramcaption{，图示取 $Q>0$。圆环每转一周，总电荷 $Q$ 就完整通过任一固定截面一次。}\label{ch21:fig:s36-5-12}
\end{minipage}\par\medskip

\begin{xesolution}
圆环转动周期为
\[
  T=\frac{2\pi}{\omega}.
\]
在一个周期内，环上全部电荷恰好通过任意固定的径向截面一次。由电流定义
\[
  \boxed{I=\frac QT=\frac{Q\omega}{2\pi}}.
\]
若 $Q>0$，电流方向与转动方向相同；若 $Q<0$，常规电流方向相反。半径 $R$ 不出现在结果中，因为
半径增大虽然使线速度增大，却同时使一周路程同比增大，周期仍只由 $\omega$ 决定。
\end{xesolution}



\begin{origexamplebox}{【真空二极管】}\label{ch21:ex:s36-5-4}
真空二极管由相对放置的阴极和阳极平板构成，两极板面积均为 $S$，距离为 $L$。稳定工作时有电流
$I$ 通过。阴极发射充分，工作于零初速的空间电荷限制状态，忽略边缘效应。取阴极处 $x=0$ 且电势为零，阳极处 $x=L$、电势为 $U>0$。依次回答：

(1) 写出两极间电势分布 $\varphi(x)$ 满足的方程；

(2) 假设电子从阴极由静止出发，求电势为 $\varphi(x)$ 处的电子速度；

(3) 利用稳恒电流与位置无关，写出电子数密度 $n(x)$；

(4) 消去 $n(x),v(x)$，得到关于 $\varphi(x)$ 的微分方程及边界条件；

(5) 积分证明查尔德--朗缪尔定律 $I=KU^{3/2}$，并求 $K$。
\par\noindent\begin{minipage}{\linewidth}
\vacuumDiodeFigure
\diagramcaption{电子由阴极射向阳极，图中 $I$ 表示相反方向的常规电流大小。}\label{ch21:fig:s36-5-13}
\end{minipage}\par\medskip
\end{origexamplebox}\par\nopagebreak[4]


\begin{xesolution}
电子云带负电，会改变板间电势，这正是空间电荷限制电流的来源。下面把 $I$ 作为电流大小。

\textbf{(1) 泊松方程。} 电子电荷密度 $\rho_q=-en(x)$，一维泊松方程为
\[
  \frac{\mathrm d^2\varphi}{\mathrm dx^2}
  =-\frac{\rho_q}{\varepsilon_0}=\frac{en(x)}{\varepsilon_0}.
\]

\textbf{(2) 能量关系。} 电子从阴极静止出发，电势能减少量转化为动能：
\[
  \frac12mv^2=e\varphi,
  \qquad v(x)=\sqrt{\frac{2e\varphi(x)}m}.
\]

\textbf{(3) 稳恒电流。} 任意截面上单位时间通过的电荷量相同，故
\[
  I=en(x)v(x)S,
  \qquad n(x)=\frac{I}{eSv(x)}.
\]

\textbf{(4) 消元与边界条件。} 代回泊松方程：
\[
  \frac{\mathrm d^2\varphi}{\mathrm dx^2}
  =A\varphi^{-1/2},
  \qquad
  A=\frac{I}{\varepsilon_0S}\sqrt{\frac m{2e}}.
\]
两板电势给出 $\varphi(0)=0,\ \varphi(L)=U$。空间电荷限制的极限状态下，阴极表面电场刚好为零，
所以还应取 $\varphi'(0)=0$。注意 $E_x=-\varphi'$，电子受力 $F_x=-eE_x=e\varphi'$：
若 $E_x(0)>0$，零初速电子受到朝阴极的力，不能向阳极出发；若 $E_x(0)<0$，电子可被拉向阳极，
但在发射不足时电流尚未达到空间电荷限制值。发射充分时，负空间电荷削弱阴极附近的加速场，
极限恰为 $E_x(0)=0$。因此这里的零场边界来自题设工作状态，不是任意稳恒电流的通用条件。

\textbf{(5) 积分。} 方程两边乘 $\varphi'$：
\[
  \varphi''\varphi'=A\varphi^{-1/2}\varphi'.
\]
从阴极积分并用 $\varphi(0)=\varphi'(0)=0$，得到
\[
  \frac12(\varphi')^2=2A\sqrt\varphi,
  \qquad \varphi'=2\sqrt A\,\varphi^{1/4}.
\]
再积分一次：
\[
  \frac43\varphi^{3/4}=2\sqrt A\,x,
  \qquad
  \varphi(x)=\left(\frac32\sqrt A\,x\right)^{4/3}.
\]
用 $\varphi(L)=U$ 消去 $A$，得到
\[
  \boxed{\varphi(x)=U\left(\frac xL\right)^{4/3}}.
\]

\par\noindent\begin{minipage}{\linewidth}
\childLangmuirFigure
\diagramcaption{解图：空间电荷限制状态下的电势曲线。阴极处切线水平，随后按 $x^{4/3}$ 增长。}
\end{minipage}\par\medskip

最后将 $A=(4/9)U^{3/2}/L^2$ 与其定义比较：
\[
  \boxed{I=\frac49\varepsilon_0S\sqrt{\frac{2e}{m}}
  \frac{U^{3/2}}{L^2}=KU^{3/2}},
  \qquad
  \boxed{K=\frac49\frac{\varepsilon_0S}{L^2}\sqrt{\frac{2e}{m}}}.
\]
\end{xesolution}



\section{伏安特性、电阻与测量}

\subsection{元件的伏安特性}


电路由元件按一定方式连接而成。对电路设计者而言，常可暂时不追究元件内部结构，而只研究端口处可测量
的电压和电流。最简单的是二端元件。图中先人为规定端口电压 $V$ 的正负方向和电流 $I$ 的测试方向；
这种规定不等于元件本身一定具有固定正、负极，计算出的负号只说明真实方向与规定方向相反。

\par\noindent\begin{minipage}{\linewidth}
\ivCharacteristicFigure
\diagramcaption{二端元件及其 $I$--$V$ 特性。标出的正、负号与箭头只是统一的测试约定。}\label{ch21:fig:s36-5-1}
\end{minipage}\par\medskip

若元件稳定工作时，同一个端电压总对应同一个端电流，就可用函数 $I=I(V)$ 描述它的外部性能，称为
伏安关系；把它画在 $I$--$V$ 平面上，便得到伏安特性曲线。不同元件的曲线可以完全不同，线性电阻
只是其中最简单的一类。使用这种二端描述还隐含两个条件：一是所研究的回路区域内没有随时间变化的磁通，
否则端电压会受感应电场影响而与路径有关；二是元件内部不能持续积累净电荷，否则流入与流出的电流并不
相等。满足这些条件后，元件的伏安关系加上连接规律就足以决定整个稳恒电路的行为。


% EM teaching enrichment 20260927
\paragraph{工作点与两种电阻}
电动势 $\mathcal E$、内阻 $r$ 的电源驱动非线性元件时，电源端口满足 $U=\mathcal E-rI$。它与元件曲线 $I=I(U)$ 的交点才是实际工作点，不能把电源电动势直接当作元件端电压。
在该点，$R_{\rm dc}=U/I$ 是静态电阻，$r_d=\dd U/\dd I$ 是小变化对应的微分电阻。若元件满足 $U=aI+bI^3$，则两者分别为 $a+bI^2$ 与 $a+3bI^2$；它们仅在线性情形相同。微分关系只能描述工作点附近的小扰动。


\subsection{欧姆定律}


欧姆定律给出一类导体两端电压与其中电流的实验关系。它并不直接解释导电的微观机制，却在广泛的材料
和工作条件下极其准确，因而成为电路分析中使用最频繁的规律之一。格奥尔格·西蒙·欧姆在十九世纪系统
研究了这一关系；后来以他的姓氏命名电阻单位欧姆，并用纪念邮票表现电阻和公式。

\ohmHistoryFigure

在温度和其他外部条件不变时，若导体两端电压为 $U$、通过电流为 $I$，欧姆定律写成
\begin{equation}
  \boxed{U=IR,\qquad I=\frac UR}.
\end{equation}
比例系数 $R$ 是电阻，单位为欧姆，$1\,\Omega=1\,\mathrm V/\mathrm A$。对长度为 $l$、横截面积为
$A$ 的均匀圆柱形导体，实验进一步给出
\begin{equation}
  \boxed{R=\rho\frac lA},
\end{equation}
其中 $\rho$ 是材料在给定状态下的电阻率。长度越长，载流子沿途经历的散射越多；截面积越大，可并行
输运电荷的通道越多，因而分别出现 $R\propto l$ 与 $R\propto1/A$。

\par\noindent\begin{minipage}{\linewidth}
\cylindricalResistorFigure
\diagramcaption{长度为 $l$、截面积为 $A$ 的均匀圆柱形电阻。电流方向与端电压测试方向一致。}\label{ch21:fig:s36-5-4}
\end{minipage}\par\medskip

理想电阻的伏安关系是一条过原点的直线。若纵轴为 $I$、横轴为 $V$，直线斜率为 $1/R$；曲线越陡，
相同电压产生的电流越大，电阻越小。

\par\noindent\begin{minipage}{\linewidth}
\ohmicIVFigure
\diagramcaption{理想电阻的 $I$--$V$ 特性。直线斜率是电导 $1/R$。}\label{ch21:fig:s36-5-5}
\end{minipage}\par\medskip

把宏观欧姆定律写成局部形式。对图~\ref{ch21:fig:s36-5-4} 的均匀导体，内部电场 $E=U/l$，电流密度 $j=I/A$。
由 $U=IR$ 与 $R=\rho l/A$ 得
\begin{equation}
  \boxed{\boldsymbol j=\frac1\rho\boldsymbol E=\sigma\boldsymbol E},
\end{equation}
其中 $\sigma=1/\rho$ 为电导率。它不含导体的长度和截面积，直接描述材料局部的导电能力。上式只在
各向同性、线性且工作状态不改变的欧姆材料中成立；对各向异性材料，$\sigma$ 一般应写成张量。


% EM teaching enrichment 20260927
\paragraph{非均匀导体的电阻积分与功率检验}
缓慢变截面的细导体可近似看成许多串接薄片，故
$\dd R=\rho(x)\dd x/A(x)$，$R=\int\rho(x)\dd x/A(x)$。若长度 $l$ 内截面积由 $A_0$ 线性增至 $2A_0$、电阻率不变，得到
\[
 R=\frac{\rho l}{A_0}\int_0^1\frac{\dd u}{1+u}
 =\frac{\rho l}{A_0}\ln2.
\]
这里采用一维电流近似；骤变截面附近的三维电流拥挤不能忽略时，应另求电流场。
局部耗散功率密度为 $\boldsymbol j\cdot\boldsymbol E=j^2/\sigma$，对体积积分应回到 $UI=I^2R$。单位和总功率是检查非均匀导体积分的有效方法。


% EM structure and teaching revision 20260928
\subsection{测量加载与电功率的局部核对}
理想电压表不取电流，理想电流表不产生压降；实际仪表会改变被测状态。电阻 $R_x$ 与内阻 $R_V$ 的电压表并联时，若电流表测的是总电流，则
\[
 R_{\text{读}}=\frac UI=\frac{R_xR_V}{R_x+R_V}<R_x.
\]
相反，若电压表跨接 $R_x$ 与电流表内阻 $R_A$ 的串联组合，而电流表只测流过 $R_x$ 的电流，则 $R_{\text{读}}=R_x+R_A$。先画出仪表端点，再分析电流究竟流过哪些支路，误差方向便不需死记。
场描述给出局部功率密度 $p=\boldsymbol j\cdot\boldsymbol E$。对欧姆介质，$p=j^2/\sigma=\sigma E^2\ge0$；积分后与端口功率 $UI$ 相等。以长直均匀电阻为例，$j=I/S$、$E=U/\ell$，故 $\int p\dd V=UI$。
这里的非负性只适用于耗散电阻。电源支路可以输出功率，其端口 $UI$ 的符号取决于是否采用电流流入电压正端的关联参考方向。


\subsection{专题例题与方法应用}

\begin{origexamplebox}{【导电厚球壳】}\label{ch21:ex:s36-5-5}
两个同心金属球壳的内、外半径分别为 $a,b$，其间充满一种特殊导电材料。该材料在各点的电导率
与当地电场大小成正比，即 $\sigma=KE$。若两球壳之间的电势差为 $U$，求两壳之间的电流。
\par\noindent\begin{minipage}{\linewidth}
\emConductingShellQuestion
\figcap{导电厚球壳的内外半径；两球面之间加给定电压。}\label{ch21:fig:em-source-thirty-sixth-section-5}
\end{minipage}\par\medskip
\end{origexamplebox}\par\nopagebreak[4]

\par\noindent\begin{minipage}{\linewidth}
\conductingShellFigure
\diagramcaption{球对称性保证电流密度和电场均沿径向，且通过任一同心球面的电流相同。}\label{ch21:fig:s36-5-14}
\end{minipage}\par\medskip

\begin{xesolution}
稳恒状态下，任意半径 $r$ 的同心球面都通过同一电流 $I$。因此
\[
  j(r)=\frac{I}{4\pi r^2}.
\]
材料的局部伏安关系不是线性的，而是
\[
  j=\sigma E=KE^2.
\]
两式联立：
\[
  E(r)=\sqrt{\frac{I}{4\pi K}}\frac1r.
\]
内外球壳的电势差为电场沿径向的积分：
\[
  U=\int_a^bE(r)\,\mathrm dr
  =\sqrt{\frac{I}{4\pi K}}\ln\frac ba.
\]
于是
\[
  \boxed{I=\frac{4\pi K}{[\ln(b/a)]^2}U^2}.
\]
这里 $I\propto U^2$，明确说明该材料不满足线性欧姆定律。若交换两球壳电势高低，电流方向反向；
上式给出其大小。
\end{xesolution}



\begin{origexamplebox}{【电离计数器】}\label{ch21:ex:s36-5-6}
最简单的电离室由一对面积为 $S$、间距为 $D$ 的平行板电极构成，置于充有混合气体的气密容器中。
阴极电势为 $-U_0$，阳极通过负载电阻 $R$ 接地。带电粒子平行极板入射，在距阳极 $x_0$ 处沿径迹
产生 $N$ 对电子--正离子；产生一对粒子所需平均能量 $W$ 由气体决定。假设全部电荷均被收集，不发生
次级电离和电子俘获，电子、正离子漂移速率分别为 $v^-,v^+$。求：

(1) 两极间电场强度；(2) 电离室电容 $C$；(3) $R$ 很大时信号电压 $\Delta U(t)$；
(4) 当 $v^-\gg v^+$ 时信号的特征；(5) $R$ 的作用不可忽略时的 $\Delta U(t)$。
\par\noindent\begin{minipage}{\linewidth}
\emIonChamberQuestion
\figcap{平行板电离室的装置、入射径迹与外接电路；竖直侧壁表示绝缘容器边界。}\label{ch21:fig:em-source-thirty-sixth-section-6}
\end{minipage}\par\medskip
\end{origexamplebox}\par\pagebreak[1]

\par\noindent\begin{minipage}{\linewidth}
\ionChamberFigure
\diagramcaption{题图：平行板电离室。$x$ 从阳极指向阴极；电子向阳极、正离子向阴极漂移并共同形成感应信号。}
\end{minipage}\par\medskip

\begin{xesolution}
若入射粒子在气体中沉积能量 $E_{\rm dep}$，则理想情况下 $N=E_{\rm dep}/W$；本题已直接给出 $N$。
以下取阳极对地电压变化为 $\Delta U$，并采用小信号条件 $Ne/C\ll U_0$，使收集过程中的场强和漂移速率
可近似不变。信号大小记为 $u=-\Delta U\ge0$，避免把实际负脉冲误画成正电压。

\textbf{(1) 电场。} 忽略边缘效应，板间电场均匀，大小
\[
  \boxed{E=\frac{U_0}{D}},
\]
方向由阳极指向阴极。电子受力方向与电场相反，故移向阳极；正离子移向阴极。

\textbf{(2) 电容。} 两平行板构成电容器：
\[
  \boxed{C=\frac{\varepsilon_0S}{D}}.
\]

\textbf{(3) $R$ 很大。} 先把两板电势保持不变。位于 $x$ 的电荷 $q$ 在阳极上感应的总电荷为
$Q_A=-q(1-x/D)$。可由叠加与静电互易关系得到：令阳极为单位电势、阴极为零，板间权重电势为
$\psi_A(x)=1-x/D$，于是 $Q_A+q\psi_A(x)=0$。也可将很薄的均匀面电荷放在两板间，
联立两侧场强跳变和总电压不变条件，得到同一关系。
电子向 $-x$ 移动、正离子向 $+x$ 移动时，保持阳极接地所需的感应正电荷都增加，增量为
$e|\Delta x|/D$。若阳极近似隔离，这部分正电荷不能由地补入，阳极电压因而降低。
在时刻 $t$，电子已经移动
$\min(v^-t,x_0)$，正离子已经移动 $\min(v^+t,D-x_0)$，所以总感应电荷大小
\[
  Q_{\rm sig}(t)=\frac{Ne}{D}
  \left[\min(v^-t,x_0)+\min(v^+t,D-x_0)\right].
\]
$R$ 极大时可忽略泄放，信号电压大小为
\[
  \boxed{|\Delta U(t)|=\frac{Ne}{CD}
  \left[\min(v^-t,x_0)+\min(v^+t,D-x_0)\right]}.
\]
最终两类电荷都到达电极，括号内等于 $D$，故 $|\Delta U(\infty)|=Ne/C$，与产生位置 $x_0$ 无关。

\textbf{(4) 快、慢两段。} 定义电子收集时间和离子收集时间
\[
  t_e=\frac{x_0}{v^-},\qquad t_i=\frac{D-x_0}{v^+}.
\]
在常见的 $t_e<t_i$ 情形，信号大小先以较大斜率快速上升；电子收集完成后，只剩离子造成的缓慢上升，
在 $t_i$ 达到平台。精确条件是 $x_0<Dv^-/(v^-+v^+)$；仅有 $v^-\gg v^+$ 不能排除极靠近阴极的
$t_i<t_e$ 情形。后一情形先失去离子分量，剩下电子分量直至 $t_e$。两种情况都由上式的两个最小值统一描述，
实际电压为 $\Delta U=-u$，最终平台为 $-Ne/C$。

\par\noindent\begin{minipage}{\linewidth}
\ionSignalFigure
\diagramcaption{解图：忽略泄放时的信号大小。图中把常见的 $t_e<t_i$ 情形画成快速前沿加缓慢离子尾。}
\end{minipage}\par\medskip

\textbf{(5) 有限负载电阻。} 读出节点相当于电容 $C$ 与电阻 $R$ 并联，信号满足
\[
  C\frac{\mathrm d u}{\mathrm dt}+\frac{u}{R}=i_s(t),\qquad \Delta U=-u.
\]
这里 $i_s$ 表示使阳极电势降低的信号电流大小，是两个矩形脉冲之和（$t\ge0$）：
\[
  i_s(t)=\frac{Ne}{D}\left[v^-H(t_e-t)+v^+H(t_i-t)\right],
\]
其中 $H$ 为阶跃函数。令 $\tau_R=RC$。对于某个持续时间为 $T_k$、幅值为
$I_k=Ne v_k/D$ 的分量，零初值解为
\[
  u_k(t)=
  \begin{cases}
    RI_k\left(1-e^{-t/\tau_R}\right),&0\le t\le T_k,\\[2mm]
    RI_k\left(1-e^{-T_k/\tau_R}\right)e^{-(t-T_k)/\tau_R},&t>T_k.
  \end{cases}
\]
分别取 $(v_k,T_k)=(v^-,t_e)$ 与 $(v^+,t_i)$，得到 $\Delta U(t)=-u_e(t)-u_i(t)$。当
$RC\gg\max(t_e,t_i)$ 时，在收集时间内展开指数，退化为前面近似无泄放的折线；当 $RC$ 较小时，信号在收集过程中
已明显经电阻衰减，收集结束后则按 $e^{-t/(RC)}$ 回到零。
\end{xesolution}



\section{电路守恒与节点方程}

\subsection{简单电阻网络}


电阻最基本的连接方式是串联和并联。串联支路中各电阻通过同一电流 $I$，总电压是各段电压之和；
并联支路中各电阻接在同一对节点之间，端电压相同，总电流是各支路电流之和。

\par\noindent\begin{minipage}{\linewidth}
\seriesParallelResistorsFigure
\diagramcaption{电阻的串联与并联。判断连接方式的依据是节点关系，而不是图形看起来是否横排或竖排。}\label{ch21:fig:s36-5-8}
\end{minipage}\par\medskip

对 $n$ 个串联电阻，$U_i=IR_i$，所以
\[
  U=U_1+\cdots+U_n=I(R_1+\cdots+R_n).
\]
等效电阻为
\begin{equation}
  \boxed{R_{\rm s}=R_1+R_2+\cdots+R_n}.
\end{equation}
对 $n$ 个并联电阻，公共端电压为 $U$，第 $i$ 条支路电流为
\begin{equation}
  \boxed{I_i=\frac{U}{R_i}}.
\end{equation}
总电流
$I=\sum_i I_i=U\sum_i(1/R_i)$，故等效电阻满足
\begin{equation}
  \boxed{\frac1{R_{\rm p}}=\frac1{R_1}+\frac1{R_2}+\cdots+\frac1{R_n}}.
\end{equation}
串联等效电阻大于任一分电阻，并联等效电阻小于任一分电阻。若网络能够逐级辨认出纯串联或纯并联，
即可反复用上述串并联公式化简；但含桥路、多个电源或交叉耦合支路的网络通常不能这样处理。


% EM teaching enrichment 20260927
\paragraph{分压、分流及仪表加载}
两电阻串联时 $U_2=UR_2/(R_1+R_2)$；两电阻并联、总电流为 $I$ 时 $I_1=IR_2/(R_1+R_2)$。分流分子是另一支路的电阻，来源是先求共同电压。
若用内阻 $R_V$ 的电压表测 $R_2$，被测支路已变为 $R_2\parallel R_V$，读数是
\[
 U_V=U\frac{R_2R_V/(R_2+R_V)}{R_1+R_2R_V/(R_2+R_V)}.
\]
只有 $R_V\gg R_2$ 且加载误差可忽略时才恢复理想分压。先把仪表画成实际元件，再讨论它的极限，可防止“测量不影响电路”的未经检验假设。


\subsection{电阻网络定理与基本术语}


一般电路是由元件按拓扑关系连接成的网络。图中的每个方框都代表一个元件，它可以是电阻、电源或更复杂
的二端器件。两个或更多元件的连接点称为节点；同一段理想导线上的各点电势相同，属于同一个节点。
相邻两节点之间包含元件的一段称为支路。若一个节点只连接两条支路，通常只是串接中的平凡节点；三条及
以上支路相接时会发生电流分流。沿若干支路从一个节点出发，最后回到原节点的闭合路径称为回路。

\par\noindent\begin{minipage}{\linewidth}
\networkTermsFigure
\diagramcaption{电路分析中的元件、节点、支路和回路。虚线只标示所选回路，不是导线。}\label{ch21:fig:s36-5-9}
\end{minipage}\par\medskip


% EM teaching enrichment 20260927
\paragraph{画图之前先识别节点}
所有用无电阻导线直接连通的点标为同一节点；交叉但未连接的导线不合并。随后把每个元件记成“节点 $a$ 到节点 $b$ 的支路”，图形折叠或拉伸便不会改变电路。
含 $N$ 个节点、$b$ 条支路的连通网络中，选 $N-1$ 条无闭环支路可连接全部节点，称为生成树。每补入一条其余支路，恰增加一个独立回路，因此独立回路数为 $b-N+1$。这是方程数量的来源，也能用来检查是否把交叉线误当节点。


\subsection{基尔霍夫电流定律（KCL）}


基尔霍夫电流定律是电荷守恒在集中参数电路中的表达：任一节点流入电流之和等于流出电流之和；
等价地，按统一符号约定计入后，节点各支路电流的代数和为零。未知支路电流的方向可以预先任意假设，
若解得负值，说明真实方向与假设方向相反。

\par\noindent\begin{minipage}{\linewidth}
\kclKvlFigure
\diagramcaption{左图为 KCL 的支路电流约定，右图为同一网络的支路电压和回路方向。所有字母标注均不覆盖元件。}\label{ch21:fig:s36-5-10}
\end{minipage}\par\medskip

按左图箭头，对四个节点分别可写为
\[
\begin{array}{ll}
N_1:&0=-i_1-i_2-i_3,\\
N_2:&0=i_1+i_4-i_6,\\
N_3:&0=i_2-i_4-i_5,\\
N_4:&0=i_3+i_5+i_6.
\end{array}
\]
最后一式中的 $i_6$ 由底边支路方向确定。把四式相加得到恒等式 $0=0$，因此含 $N$ 个节点的连通网络
只有 $N-1$ 个独立 KCL 方程。少写会丢失约束，多写则只是重复已有信息。


% EM teaching enrichment 20260927
\paragraph{节点暂时储存电荷时的完整守恒式}
用闭面包围一个小节点，流入减流出等于 $\dd Q_{\rm node}/\dd t$。稳恒状态右边为零，才得到常用 KCL。
在集中参数模型中，若节点对地有不可忽略电容 $C$，就应把 $C\dd V/\dd t$ 作为一条支路电流列入。例如通过电阻 $R$ 给该节点充电时，
\[
 \frac{\mathcal E-V}{R}=C\frac{\dd V}{\dd t}.
\]
常用节点法不是否定电荷积累，而是把积累明确归入电容元件。先选电流正向再写等式，解出的负数不应被擅自改成正数。


\subsection{基尔霍夫电压定律（KVL）}


基尔霍夫电压定律是稳恒电路能量守恒的等价表达：沿任一闭合回路，各支路电压的代数和为零。使用时
先规定每个支路电压的正方向，再选定绕行方向；沿支路从正端走向负端记为一个符号，反向走则记为相反符号。
对图~\ref{ch21:fig:s36-5-10} 右侧标出的四个回路，可写成
\[
\begin{array}{ll}
L_1:&0=-v_1+v_2+v_4,\\
L_2:&0=-v_2+v_3-v_5,\\
L_3:&0=-v_4+v_5-v_6,\\
L_4:&0=v_1+v_6-v_3.
\end{array}
\]
四式相加同样得到恒等式，所以它们并不独立。前三个小回路已经覆盖全部支路，外围回路 $L_4$ 可由前三式
线性组合得到。对连通网络，独立回路数等于“支路数 $-$ 节点数 $+1$”。KCL、KVL 再配合每个元件的
伏安关系，原则上可以求出全部支路电流和元件电压。


% EM teaching enrichment 20260927
\paragraph{一圈电压与一圈能量如何对应}
对电动势 $\mathcal E$、内阻 $r$、外电阻 $R$ 的闭合回路，沿电流绕行，源内非静电作用提升每单位电荷的能量 $\mathcal E$，电阻段合计降低 $I(R+r)$，故
$\mathcal E-Ir-IR=0$。
乘以 $I$ 得电源输出功率 $\mathcal EI=I^2r+I^2R$，同时核对了符号与能量。
若回路有变化磁通，应加入感应电动势；不能一面把感应效应遗漏，一面仍以纯静电势差的闭合和为零处理全部电场。


\subsection{节点电压法}


直接把所有支路电流和元件电压都列为未知量，哪怕小网络也会产生很大的方程组。节点电压法先选一个节点
作为参考地，把其电势定为零；随后只把其余独立节点电势作为未知量。支路电压由两端节点电势之差自动给出，
因此 KVL 自动满足；再对未知节点写 KCL，就得到数量最少的一组方程。

\par\noindent\begin{minipage}{\linewidth}
\nodeMethodFigure
\diagramcaption{节点电压法示例。取 $c$ 为参考地，理想电压源固定 $a$ 点电势，只剩 $b$ 点电势未知。}\label{ch21:fig:s36-5-11}
\end{minipage}\par\medskip

在图中令 $\varphi_c=0$，电压源使 $\varphi_a=E_0$，设 $\varphi_b=e$。以流入 $b$ 点为正：
从 $a$ 经 $R_1$ 流入的电流是 $(E_0-e)/R_1$，电流源注入 $I_0$，从 $b$ 经 $R_2$ 流出的电流是
$e/R_2$。KCL 给出
\[
  \frac{E_0-e}{R_1}+I_0-\frac e{R_2}=0.
\]
解得
\[
  \boxed{e=\frac{E_0/R_1+I_0}{1/R_1+1/R_2}}.
\]
节点电势一旦求出，每个元件两端的电压和各支路电流便可由其伏安关系依次得到。参考节点原则上可以任意
选择；经验上，选与较多理想电压源相连的节点作参考地，往往能减少辅助未知量，使代数最简。

% EM teaching enrichment 20260927
\paragraph{悬浮电压源与解后的核对}
若理想电压源位于两个未知节点 $a,b$ 之间，不能把其电流写成“电压除以零电阻”。可把两节点合成超节点，对流向外界的支路写 KCL，再补 $V_a-V_b=\mathcal E$。
例如 $a,b$ 分别经 $R_a,R_b$ 接地，则
$V_a/R_a+V_b/R_b=0$，与电压约束联立得
$V_a=\mathcal E R_a/(R_a+R_b)$、$V_b=-\mathcal E R_b/(R_a+R_b)$。
负节点电势是相对参考点而言，并非异常。解后应检查各节点电流守恒，以及全部元件吸收功率的代数和为零。




% EM structure and teaching revision 20260928
\subsection{把拓扑、方程与功率检查分开}
设连通电路有 $N$ 个节点和 $B$ 条支路。固定一个参考电势后，可独立写出 $N-1$ 个节点电流方程；若改用支路电流未知量，则还需 $B-N+1$ 个独立回路方程及元件关系。逐个面孔写回路虽方便，但非平面画法中的线条交叉未必表示新节点。
例如节点电势 $V$ 经 $R_1$ 接到已知电势 $E$、经 $R_2$ 接地，同时有恒流源 $I_s$ 向该节点注入，则
\[
 \frac{V-E}{R_1}+\frac V{R_2}-I_s=0,
 \qquad V=\frac{E/R_1+I_s}{1/R_1+1/R_2}.
\]
令 $I_s=0$ 应退回分压公式；令 $R_2\to\infty$ 应得到 $V=E+I_sR_1$。这类极限检查比单纯把计算结果代回原方程更容易发现参考方向选错。
若理想电压源连接两个未知电势节点，把两节点作为超节点写总电流平衡，再补电压差约束。解完后分别计算各电阻消耗和各电源输出，检查全电路净功率为零；这一步同时检验数值、支路方向和正负号。


\subsection{专题例题与方法应用}

\begin{origexamplebox}{【圆环电阻】}\label{ch21:ex:s36-5-7}
一段总电阻为 $R$ 的均匀电阻丝弯成圆环。在环上固定一点 $a$，另一接点 $b$ 可沿电阻丝移动。
先在 $a,b$ 间接恒定电压 $U$，问 $b$ 在何处时整条电阻丝的发热功率最小，最小值是多少？若改接
电流为 $I$ 的恒流源，结论又如何？
\par\noindent\begin{minipage}{\linewidth}
\emRingResistorQuestion
\figcap{均匀电阻圆环及接点；两段弧的电阻记为 $R_1,R_2$。}\label{ch21:fig:em-source-thirty-sixth-section-7}
\end{minipage}\par\medskip
\end{origexamplebox}\par\nopagebreak[4]

\par\noindent\begin{minipage}{\linewidth}
\ringResistorFigure
\diagramcaption{题图：接点 $a,b$ 把圆环分成电阻为 $R_1,R_2$ 的两段，两段端点相同，故彼此并联。}
\end{minipage}\par\medskip

\begin{xesolution}
两段弧的电阻之和固定：$R_1+R_2=R$。它们并联后的等效电阻为
\[
  R_{\rm eq}=\frac{R_1R_2}{R_1+R_2}=\frac{R_1R_2}{R}.
\]

\textbf{恒压源。} 总功率
\[
  P_U=\frac{U^2}{R_{\rm eq}}=\frac{U^2R}{R_1R_2}.
\]
在和固定时，乘积满足 $R_1R_2\le(R/2)^2$；要使功率最小，应使乘积最大，故
\[
  R_1=R_2=\frac R2.
\]
均匀电阻丝上这意味着 $b$ 位于 $a$ 的直径对面，且
\[
  \boxed{P_{U,\min}=\frac{4U^2}{R}}.
\]

\textbf{恒流源。} 此时
\[
  P_I=I^2R_{\rm eq}=I^2\frac{R_1R_2}{R}.
\]
要使功率最小，应使 $R_1R_2$ 尽可能小，即让 $b$ 趋近 $a$，一段电阻趋近零，因而
\[
  \boxed{P_{I,\min}=0}
\]
（若规定 $a,b$ 必须是不同点，则 $0$ 是下确界，只能无限接近而不能取到）。相反，$b$ 位于直径对面
时 $R_{\rm eq}$ 最大，恒流情况下反而得到最大功率 $I^2R/4$。这两种结论相反，根源是恒压时
$P=U^2/R_{\rm eq}$，恒流时 $P=I^2R_{\rm eq}$。
\end{xesolution}



\begin{origexamplebox}{【双电源供电】}\label{ch21:ex:s36-5-8}
两个电动势和内阻分别为 $(\varepsilon_1,r_1)$、$(\varepsilon_2,r_2)$ 的实际电源同极性并联，
共同连接负载电阻 $R$。

(1) 利用基尔霍夫方程求通过 $R$ 的电流；

(2) 对结果变形，说明它怎样给出关于电源叠加和戴维南等效的启发。
\par\noindent\begin{minipage}{\linewidth}
\dualSourceFigure
\diagramcaption{每个实际电源均由理想电动势与内阻串联构成，三条支路跨接同一对节点。}\label{ch21:fig:s36-5-15}
\end{minipage}\par\medskip
\end{origexamplebox}\par\nopagebreak[4]


\begin{xesolution}
记左、右公共节点电势差为 $V=V_L-V_R$，并把三条支路电流都暂按从左向右取正。负载支路电流为
$V/R$；对第 $k$ 个电源支路，按图示极性有 $V=\varepsilon_k+i_kr_k$，所以
$i_k=(V-\varepsilon_k)/r_k$。在左节点写 KCL：
\[
  \frac VR+\frac{V-\varepsilon_1}{r_1}
  +\frac{V-\varepsilon_2}{r_2}=0.
\]
解出 $V$ 并除以 $R$，负载电流为
\[
  \boxed{I=\frac VR
  =\frac{\varepsilon_1r_2+\varepsilon_2r_1}
  {r_1r_2+R(r_1+r_2)}}.
\]

把两个电源及其内阻看成一个二端网络。定义
\[
  \varepsilon_{\rm eq}
  =\frac{\varepsilon_1/r_1+\varepsilon_2/r_2}{1/r_1+1/r_2}
  =\frac{\varepsilon_1r_2+\varepsilon_2r_1}{r_1+r_2},
  \qquad
  r_{\rm eq}=\frac{r_1r_2}{r_1+r_2}.
\]
原结果立即化成
\[
  \boxed{I=\frac{\varepsilon_{\rm eq}}{R+r_{\rm eq}}}.
\]
这说明并联电源网络从负载端看，等效为一个电动势 $\varepsilon_{\rm eq}$ 与内阻 $r_{\rm eq}$ 串联的
实际电源，正是戴维南等效的形式。同时，$V$ 对 $\varepsilon_1,\varepsilon_2$ 是线性的，每个电源的
贡献可以分别计算再相加，体现了线性网络的叠加思想。若两电源极性相反，只需把其中一个
$\varepsilon_k$ 按符号取负；公式仍然成立。
\end{xesolution}
}

{
\graphicspath{{/Users/dongxingxi/Documents/物理竞赛一轮资料/sample/xes-original/thirty-seventh-section/}}
\section{等效电源、对称性与电桥}

% Technical figures are redrawn as vectors.  Every text node is transparent;
% rectangles are used only for actual circuit elements or explanatory regions.
\tikzset{
  xeslabel/.style={font=\small,text=Ink,fill=none,inner sep=.45pt},
  xessmall/.style={font=\footnotesize,text=TextGray,fill=none,inner sep=.35pt},
  xestiny/.style={font=\scriptsize,text=TextGray,fill=none,inner sep=.25pt},
  circuit/.style={draw=Ink,line width=.82pt,line cap=round,line join=round},
  activewire/.style={draw=DeepBlue,line width=1.05pt,line cap=round,line join=round},
  hiddenwire/.style={draw=TextGray,densely dashed,line width=.72pt,line cap=round},
  solutionwire/.style={draw=DeepBlue!78,line width=.9pt,line cap=round,line join=round},
  currentarrow/.style={-{Latex[length=1.8mm,width=1.2mm]},draw=DeepBlue,line width=.75pt},
  nodepoint/.style={circle,fill=Ink,draw=none,inner sep=1.45pt},
  bluenode/.style={circle,fill=DeepBlue,draw=white,line width=.3pt,inner sep=1.7pt},
  goldnode/.style={circle,fill=Gold!85!black,draw=white,line width=.3pt,inner sep=1.7pt},
  port/.style={circle,fill=white,draw=Ink,line width=.75pt,inner sep=1.6pt},
  sourcecircle/.style={circle,draw=Ink,line width=.75pt,minimum size=.54cm,inner sep=0pt},
  networkblob/.style={draw=DeepBlue!75,fill=DeepBlue!4,line width=.8pt},
  eqbox/.style={draw=none,fill=none,inner sep=2pt}
}


\newcommand{\zigres}[2]{%
  \draw[circuit,decorate,decoration={zigzag,segment length=2.25mm,amplitude=.72mm}] (#1)--(#2);}

\newcommand{\drawvoltage}[2]{%
  \draw[circuit] (#1)--++(0,.26);
  \draw[circuit] ($(#1)+(-.16,.26)$)--($(#1)+(.16,.26)$);
  \draw[circuit] ($(#1)+(-.27,.39)$)--($(#1)+(.27,.39)$);
  \draw[circuit] ($(#1)+(0,.39)$)--(#2);}

\newcommand{\superpositionFigure}{%
\begin{center}
\begin{tikzpicture}[x=1cm,y=1cm]
  \begin{scope}[shift={(-3.25,0)}]
    \coordinate (a) at (-1.7,1.05);\coordinate (b) at (.15,1.05);\coordinate (c) at (.15,-1.0);
    \coordinate (d) at (-1.7,-1.0);\coordinate (r) at (1.6,0);
    \draw[circuit] (a)--(-1.7,.32);
    \draw[circuit] (-1.93,.18)--(-1.47,.18);
    \draw[circuit] (-1.82,.02)--(-1.58,.02);
    \draw[circuit] (-1.7,.02)--(d);
    \zigres{a}{b}\zigres{b}{c}\draw[circuit] (c)--(d);
    \draw[circuit] (b)--(1.55,1.05)--(1.55,.55);
    \draw[circuit] (c)--(1.55,-1.0)--(1.55,-.50);
    \node[nodepoint] at (a) {};\node[nodepoint] at (b) {};\node[nodepoint] at (c) {};
    \node[xeslabel,above] at (a) {$a$};\node[xeslabel,above] at (b) {$b$};\node[xeslabel,below] at (c) {$c$};
    \node[xestiny,above] at (-.77,1.28) {$R_1$};\node[xestiny,right] at (.40,.03) {$R_2$};
    \node[xestiny,left] at (-1.98,-.02) {$E_0$};
    \node[xessmall] at (0,-1.55) {只保留电压源：电流源断路};
  \end{scope}
  \begin{scope}[shift={(3.25,0)}]
    \coordinate (a) at (-1.7,1.05);\coordinate (b) at (.15,1.05);\coordinate (c) at (.15,-1.0);
    \draw[circuit] (a)--(-1.7,-1.0)--(c);
    \zigres{a}{b}\zigres{b}{c}
    \draw[circuit] (b)--(1.45,1.05)--(1.45,.36);
    \draw[circuit] (c)--(1.45,-1.0)--(1.45,-.36);
    \node[sourcecircle] at (1.45,0) {};
    \draw[currentarrow] (1.45,-.22)--(1.45,.22);
    \node[xestiny,right] at (1.73,0) {$I_0$};
    \node[nodepoint] at (a) {};\node[nodepoint] at (b) {};\node[nodepoint] at (c) {};
    \node[xeslabel,above] at (a) {$a$};\node[xeslabel,above] at (b) {$b$};\node[xeslabel,below] at (c) {$c$};
    \node[xestiny,above] at (-.77,1.28) {$R_1$};\node[xestiny,right] at (.40,.03) {$R_2$};
    \node[xessmall] at (0,-1.55) {只保留电流源：电压源短路};
  \end{scope}
\end{tikzpicture}
\end{center}}

\newcommand{\networkBlob}[1]{%
  \draw[networkblob] plot[smooth cycle,tension=.7] coordinates
    {(-1.05,.72) (-.55,1.05) (.10,.85) (.82,1.02) (1.08,.35) (.92,-.38) (.34,-.86) (-.40,-.78) (-1.04,-.25)};
  \zigres{-0.65,0.42}{0.20,0.42}
  \node[sourcecircle,minimum size=.42cm] at (-.45,-.35) {};
  \draw[currentarrow] (-.45,-.50)--(-.45,-.20);
  \draw[circuit] (.42,-.55)--(.42,-.18);\draw[circuit] (.24,-.18)--(.60,-.18);
  \draw[circuit] (.15,-.05)--(.69,-.05);\draw[circuit] (.42,-.05)--(.42,.15);
  \coordinate (#1u) at (1.72,.34);\coordinate (#1d) at (1.72,-.34);
  \draw[circuit] (1.00,.34)--(#1u);\draw[circuit] (1.00,-.34)--(#1d);
  \node[port] at (#1u) {};\node[port] at (#1d) {};
}

\newcommand{\theveninDerivationFigure}{%
\begin{center}
\begin{tikzpicture}[x=.82cm,y=.82cm]
  \begin{scope}[shift={(-6.2,0)}]
    \networkBlob{pA}
    \draw[currentarrow] (2.72,.34)--(2.05,.34);
    \node[xestiny,above] at (2.38,.43) {$i$};
    \node[xestiny,right] at (1.92,.25) {$+$};\node[xestiny,right] at (1.92,-.43) {$-$};
    \node[xestiny] at (2.28,-.12) {$v$};\node[xessmall] at (.25,-1.43) {(a) 原有源二端网络};
  \end{scope}
  \begin{scope}[shift={(-2.05,0)}]
    \networkBlob{pB}
    \node[sourcecircle] at (2.45,0) {};
    \draw[currentarrow] (2.45,-.22)--(2.45,.22);
    \draw[circuit] (pBu)--(2.45,.34)--(2.45,.27) (pBd)--(2.45,-.34)--(2.45,-.27);
    \draw[currentarrow] (2.05,.70)--(1.42,.70);
    \node[xestiny,above] at (1.74,.78) {$i_t$};
    \node[xestiny,right] at (2.78,.02) {$i_t$};\node[xestiny,below] at (2.45,-.43) {$v_t$};
    \node[xessmall] at (.45,-1.43) {(b) 接入测试电流源};
  \end{scope}
  \begin{scope}[shift={(2.45,0)}]
    \draw[networkblob] plot[smooth cycle,tension=.7] coordinates
      {(-1.05,.72) (-.4,1.0) (.28,.83) (.9,.98) (1.08,.30) (.88,-.55) (.14,-.85) (-.72,-.65) (-1.05,-.05)};
    \zigres{-0.62,0.35}{0.28,0.35}
    \draw[circuit] (-.55,-.35)--(.56,-.35);
    \draw[circuit] (1.0,.34)--(1.72,.34) (1.0,-.34)--(1.72,-.34);
    \node[port] at (1.72,.34) {};\node[port] at (1.72,-.34) {};
    \node[sourcecircle] at (2.45,0) {};\draw[currentarrow] (2.45,-.22)--(2.45,.22);
    \draw[circuit] (1.72,.34)--(2.45,.34)--(2.45,.27) (1.72,-.34)--(2.45,-.34)--(2.45,-.27);
    \node[xestiny,right] at (2.82,.02) {$i_t$};\node[xestiny,below] at (2.45,-.43) {$v_a=i_tr$};
    \node[xessmall] at (.55,-1.43) {(c) 内部独立源置零};
  \end{scope}
  \begin{scope}[shift={(6.6,0)}]
    \networkBlob{pD}
    \draw[circuit] (pDu)--++(.48,0);\draw[circuit] (pDd)--++(.48,0);
    \node[xestiny,right] at (1.92,.25) {$+$};\node[xestiny,right] at (1.92,-.43) {$-$};
    \node[xestiny,below] at (2.10,-.48) {$v_b$};\node[xestiny,above] at (2.25,.48) {$i_t=0$};
    \node[xessmall] at (.35,-1.43) {(d) 端口开路};
  \end{scope}
\end{tikzpicture}
\end{center}}

\newcommand{\sourceEquivalenceFigure}{%
\begin{center}
\begin{tikzpicture}[x=1cm,y=1cm]
  \begin{scope}[shift={(-4.65,0)}]
    \networkBlob{qA}
    \node[xessmall] at (.30,-1.28) {线性有源网络};
  \end{scope}
  \node[xeslabel] at (-1.95,0) {$\Longleftrightarrow$};
  \begin{scope}[shift={(0,0)}]
    \coordinate (u) at (-1.1,.72);\coordinate (d) at (-1.1,-.72);
    \node[sourcecircle] at (-1.1,0) {};
    \node[xestiny] at (-1.1,.12) {$+$};\node[xestiny] at (-1.1,-.12) {$-$};
    \draw[circuit] (u)--(-1.1,.27) (d)--(-1.1,-.27);
    \draw[circuit] (u)--(-.42,.72);\zigres{-0.42,.72}{.70,.72}\draw[circuit] (.70,.72)--(1.30,.72);
    \draw[circuit] (d)--(1.30,-.72);\node[port] at (1.30,.72) {};\node[port] at (1.30,-.72) {};
    \node[xestiny,below] at (-1.10,-.34) {$V_{\rm Th}$};\node[xestiny,above] at (.14,.95) {$R_{\rm Th}$};
    \node[xessmall] at (0,-1.28) {戴维南形式};
  \end{scope}
  \node[xeslabel] at (2.15,0) {$\Longleftrightarrow$};
  \begin{scope}[shift={(4.35,0)}]
    \draw[circuit] (-1.25,.72)--(1.25,.72) (-1.25,-.72)--(1.25,-.72);
    \draw[circuit] (-.85,.72)--(-.85,.27) (-.85,-.72)--(-.85,-.27);
    \node[sourcecircle] at (-.85,0) {};\draw[currentarrow] (-.85,-.22)--(-.85,.22);
    \zigres{.25,.72}{.25,-.72}
    \node[port] at (1.25,.72) {};\node[port] at (1.25,-.72) {};
    \node[xestiny,left] at (-1.15,0) {$I_{\rm N}$};\node[xestiny,right] at (.52,0) {$R_{\rm N}$};
    \node[xessmall] at (0,-1.28) {诺顿形式};
  \end{scope}
\end{tikzpicture}
\end{center}}

\newcommand{\bridgeFigure}{%
\begin{center}
\begin{tikzpicture}[x=1cm,y=1cm]
  \coordinate (L) at (-3,0);\coordinate (T) at (0,1.25);\coordinate (D) at (0,-1.25);\coordinate (R) at (3,0);
  \zigres{L}{T}\zigres{T}{R}\zigres{L}{D}\zigres{D}{R}\zigres{T}{D}
  \draw[circuit] (L)--(-3,-2.0)--(-.55,-2.0);
  \draw[circuit] (.55,-2.0)--(3,-2.0)--(R);
  \draw[circuit] (-.55,-2.0)--(-.18,-2.0);\draw[circuit] (.18,-2.0)--(.55,-2.0);
  \draw[circuit] (-.18,-1.72)--(-.18,-2.28);\draw[circuit] (.18,-1.83)--(.18,-2.17);
  \foreach \P in {L,T,D,R}{\node[nodepoint] at (\P) {};}
  \node[xestiny,above left] at (-1.55,.70) {$R_1$};\node[xestiny,above right] at (1.55,.70) {$R_2$};
  \node[xestiny,below left] at (-1.55,-.70) {$R_3$};\node[xestiny,below right] at (1.55,-.70) {$R_4$};
  \node[xestiny,right] at (.25,0) {$R_5$};\node[xestiny,below] at (0,-2.28) {$\mathcal E$};
  \draw[currentarrow] (.43,.54)--(.43,-.54);\node[xestiny,right] at (.54,.02) {$I_5=0$};
\end{tikzpicture}
\end{center}}

\newcommand{\deltaStarFigure}{%
\begin{center}
\begin{tikzpicture}[x=1cm,y=1cm]
  \begin{scope}[shift={(-3.4,0)}]
    \coordinate (A) at (0,1.75);\coordinate (B) at (-1.75,-1.25);\coordinate (C) at (1.75,-1.25);
    \zigres{A}{B}\zigres{A}{C}\zigres{B}{C}
    \foreach \P in {A,B,C}{\node[port] at (\P) {};}
    \node[xeslabel,above] at (A) {$A$};\node[xeslabel,below left] at (B) {$B$};\node[xeslabel,below right] at (C) {$C$};
    \node[xestiny,left] at (-1.05,.38) {$R_{AB}$};\node[xestiny,right] at (1.05,.38) {$R_{AC}$};
    \node[xestiny,below] at (0,-1.48) {$R_{BC}$};\node[xessmall] at (0,-2.05) {$\Delta$ 连接};
  \end{scope}
  \node[xeslabel] at (0,0) {$\Longleftrightarrow$};
  \begin{scope}[shift={(3.4,0)}]
    \coordinate (A) at (0,1.75);\coordinate (B) at (-1.75,-1.25);\coordinate (C) at (1.75,-1.25);\coordinate (O) at (0,0);
    \zigres{A}{O}\zigres{B}{O}\zigres{C}{O}
    \foreach \P in {A,B,C}{\node[port] at (\P) {};}
    \node[nodepoint] at (O) {};
    \node[xeslabel,above] at (A) {$A$};\node[xeslabel,below left] at (B) {$B$};\node[xeslabel,below right] at (C) {$C$};
    \node[xestiny,right] at (.28,.93) {$R_a$};\node[xestiny,above left] at (-.95,-.50) {$R_b$};
    \node[xestiny,above right] at (.95,-.50) {$R_c$};\node[xestiny,above right] at (O) {$\varphi_0$};
    \node[xessmall] at (0,-2.05) {$Y$ 连接};
  \end{scope}
\end{tikzpicture}
\end{center}}

\newcommand{\sourceTransformExampleFigure}{%
\begin{center}
\begin{tikzpicture}[x=1cm,y=1cm]
  \begin{scope}[shift={(-3.2,0)}]
    \node[sourcecircle] at (-1.55,0) {};
    \node[xestiny] at (-1.55,.12) {$+$};\node[xestiny] at (-1.55,-.12) {$-$};
    \draw[circuit] (-1.55,.27)--(-1.55,.90)--(-.85,.90);
    \zigres{-0.85,.90}{.35,.90}\draw[circuit] (.35,.90)--(1.40,.90);
    \draw[circuit] (-1.55,-.27)--(-1.55,-.90)--(1.40,-.90);
    \node[port] at (1.40,.90) {};\node[port] at (1.40,-.90) {};
    \draw[currentarrow] (.92,1.24)--(1.38,1.24);
    \node[xestiny,above] at (1.15,1.28) {$I$};
    \node[xestiny,left] at (-1.88,0) {$E$};\node[xestiny,above] at (-.25,1.12) {$r$};
    \node[xestiny,right] at (1.60,.68) {$+$};\node[xestiny,right] at (1.60,-1.02) {$-$};
    \node[xestiny] at (1.85,0) {$U$};\node[xessmall] at (0,-1.45) {电压源形式};
  \end{scope}
  \node[xeslabel] at (0,0) {$\Longleftrightarrow$};
  \begin{scope}[shift={(3.2,0)}]
    \draw[circuit] (-1.45,.90)--(1.40,.90) (-1.45,-.90)--(1.40,-.90);
    \draw[circuit] (-.85,.90)--(-.85,.27) (-.85,-.90)--(-.85,-.27);
    \node[sourcecircle] at (-.85,0) {};\draw[currentarrow] (-.85,-.22)--(-.85,.22);
    \zigres{.20,.90}{.20,-.90}
    \node[port] at (1.40,.90) {};\node[port] at (1.40,-.90) {};
    \draw[currentarrow] (.92,1.24)--(1.38,1.24);
    \node[xestiny,above] at (1.15,1.28) {$I$};
    \node[xestiny,left] at (-1.17,0) {$E/r$};\node[xestiny,right] at (.47,0) {$r$};
    \node[xestiny,right] at (1.60,.68) {$+$};\node[xestiny,right] at (1.60,-1.02) {$-$};
    \node[xestiny] at (1.85,0) {$U$};\node[xessmall] at (0,-1.45) {电流源形式};
  \end{scope}
\end{tikzpicture}
\end{center}}

\newcommand{\gridProblemFigure}{%
\begin{center}
\begin{tikzpicture}[x=1.05cm,y=1.05cm]
  \foreach \i in {0,1,2}{\draw[circuit] (0,\i)--(2,\i);\draw[circuit] (\i,0)--(\i,2);}
  \foreach \i in {0,1,2}{\foreach \j in {0,1,2}{\node[nodepoint] at (\i,\j) {};}}
  \node[xeslabel,below left] at (0,0) {$A$};\node[xeslabel,below right] at (2,0) {$B$};
  \node[xeslabel,above right] at (2,2) {$C$};
  \draw[<->,DeepBlue!75,line width=.6pt] (0,-.55)--(1,-.55);
  \node[xestiny,below] at (.5,-.60) {$R$};
\end{tikzpicture}
\end{center}}

\newcommand{\potentialNode}[3]{%
  \node[circle,fill=#3,draw=white,line width=.3pt,inner sep=1.55pt] at (#1) {};
  \node[xestiny,text=Ink,anchor=south west] at ($(#1)+(.08,.08)$) {$#2$};}

\newcommand{\gridPotentialDiagonalFigure}{%
\begin{center}
\begin{tikzpicture}[x=1.2cm,y=1.2cm]
  \foreach \i in {0,1,2}{\draw[circuit] (0,\i)--(2,\i);\draw[circuit] (\i,0)--(\i,2);}
  \potentialNode{0,0}{1}{Gold!70}
  \potentialNode{1,0}{\frac23}{DeepBlue!28}
  \potentialNode{2,0}{\frac12}{RuleGray!55}
  \potentialNode{0,1}{\frac23}{DeepBlue!28}
  \potentialNode{1,1}{\frac12}{RuleGray!55}
  \potentialNode{2,1}{\frac13}{DeepBlue!12}
  \potentialNode{0,2}{\frac12}{RuleGray!55}
  \potentialNode{1,2}{\frac13}{DeepBlue!12}
  \potentialNode{2,2}{0}{Ink!20}
  \node[xestiny,below left] at (0,0) {$A$};\node[xestiny,below right=3pt] at (2,2) {$C$};
\end{tikzpicture}
\end{center}}

\newcommand{\gridPotentialAdjacentFigure}{%
\begin{center}
\begin{tikzpicture}[x=1.2cm,y=1.2cm]
  \foreach \i in {0,1,2}{\draw[circuit] (0,\i)--(2,\i);\draw[circuit] (\i,0)--(\i,2);}
  \potentialNode{0,0}{1}{Gold!70}
  \potentialNode{1,0}{\frac12}{RuleGray!55}
  \potentialNode{2,0}{0}{Ink!20}
  \potentialNode{0,1}{\frac7{10}}{DeepBlue!30}
  \potentialNode{1,1}{\frac12}{RuleGray!55}
  \potentialNode{2,1}{\frac3{10}}{DeepBlue!12}
  \potentialNode{0,2}{\frac35}{DeepBlue!23}
  \potentialNode{1,2}{\frac12}{RuleGray!55}
  \potentialNode{2,2}{\frac25}{DeepBlue!16}
  \node[xestiny,below left] at (0,0) {$A$};\node[xestiny,below right] at (2,0) {$B$};
\end{tikzpicture}
\end{center}}

\newcommand{\tetrahedronProblemFigure}{%
\begin{center}
\begin{tikzpicture}[x=1cm,y=1cm]
  \coordinate (T) at (0,2.0);\coordinate (L) at (-2.0,.15);\coordinate (R) at (2.0,.15);\coordinate (B) at (0,-1.65);
  \draw[circuit] (T)--(L)--(B)--(R)--(T) (L)--(R) (T)--(B);
  \foreach \P in {T,L,R,B}{\node[bluenode] at (\P) {};}
  \node[xestiny,above] at (T) {$A$};\node[xestiny,left] at (L) {$C$};
  \node[xestiny,right] at (R) {$D$};\node[xestiny,below] at (B) {$B$};
  \node[xessmall] at (0,-2.15) {交叉点只是投影重叠，不是电路节点};
\end{tikzpicture}
\end{center}}

\newcommand{\completeGraphProblemFigure}{%
\begin{center}
\begin{tikzpicture}[x=1cm,y=1cm]
  \foreach \i in {0,...,4}{\coordinate (P\i) at ({90+72*\i}:1.75);}
  \foreach \i in {0,...,4}{\foreach \j in {0,...,4}{\ifnum\i<\j\draw[circuit] (P\i)--(P\j);\fi}}
  \foreach \i in {0,...,4}{\node[bluenode] at (P\i) {};}
  \node[xestiny,above] at (P0) {$N=5$};
\end{tikzpicture}
\end{center}}

\newcommand{\completeGraphReductionFigure}{%
\begin{center}
\begin{tikzpicture}[x=1cm,y=1cm]
  \coordinate (A) at (-3,0);\coordinate (B) at (3,0);
  \node[goldnode] at (A) {};\node[goldnode] at (B) {};
  \node[xeslabel,left] at (A) {$A:V$};\node[xeslabel,right] at (B) {$B:0$};
  \draw[activewire] (A)--(B);
  \foreach \y in {1.25,.42,-.42,-1.25}{
    \coordinate (M) at (0,\y);\draw[solutionwire] (A)--(M)--(B);\node[bluenode] at (M) {};}
  \draw[RuleGray,densely dashed,rounded corners=3pt] (-.45,1.60) rectangle (.45,-1.60);
  \node[xessmall,anchor=west] at (.58,1.42) {$N-2$ 个等势点：$V/2$};
  \node[xestiny] at (-1.5,1.55) {$R$};\node[xestiny] at (1.5,1.55) {$R$};
  \node[xessmall] at (0,-1.95) {每条支路段均为 $R$；包括直接连接 $A,B$ 的一段};
\end{tikzpicture}
\end{center}}

\newcommand{\octagonProblemFigure}{%
\begin{center}
\begin{tikzpicture}[x=1cm,y=1cm]
  \foreach \i in {0,...,7}{\coordinate (O\i) at ({22.5+45*\i}:2.0);}
  \foreach \i in {0,...,7}{
    \foreach \j in {0,...,7}{
      \ifnum\i<\j
        \pgfmathtruncatemacro{\d}{mod(\j-\i,8)}
        \ifnum\d=4\else\draw[circuit,line width=.55pt] (O\i)--(O\j);\fi
      \fi}}
  \foreach \i in {0,...,7}{\node[bluenode] at (O\i) {};}
  \node[xestiny,right] at (O0) {$A$};\node[xestiny,left] at (O4) {$B$};
  \draw[DeepBlue!70,densely dashed,line width=.7pt] (O0)--(O4);
  \node[xessmall,below] at (0,-2.38) {蓝色虚线直径不接电阻；其余任意顶点对均有 $R$};
\end{tikzpicture}
\end{center}}

\newcommand{\octagonReductionFigure}{%
\begin{center}
\begin{tikzpicture}[x=1cm,y=1cm]
  \coordinate (A) at (-3.2,0);\coordinate (B) at (3.2,0);
  \node[goldnode] at (A) {};\node[goldnode] at (B) {};
  \node[xeslabel,left] at (A) {$A$};\node[xeslabel,right] at (B) {$B$};
  \foreach \y in {1.55,.95,.32,-.32,-.95,-1.55}{
    \coordinate (M) at (0,\y);\draw[solutionwire] (A)--(M)--(B);\node[bluenode] at (M) {};}
  \node[xestiny] at (-1.62,1.75) {$R$};\node[xestiny] at (1.62,1.75) {$R$};
  \draw[RuleGray,densely dashed,rounded corners=3pt] (-.45,1.82) rectangle (.45,-1.82);
  \node[xessmall] at (0,-2.12) {六个中间点等势};
  \draw[DeepBlue!70,densely dashed] (A)--(B);
  \node[xessmall] at (0,-2.62) {虚线只表示直径位置，$A,B$ 间无直接支路};
\end{tikzpicture}
\end{center}}

\newcommand{\cubeCoordinates}{%
  \coordinate (A) at (-1.35,-1.05);\coordinate (B) at (1.05,-1.05);
  \coordinate (E) at (-1.35,1.10);\coordinate (F) at (1.05,1.10);
  \coordinate (D) at (-.45,-.32);\coordinate (C) at (1.95,-.32);
  \coordinate (H) at (-.45,1.83);\coordinate (G) at (1.95,1.83);}

\newcommand{\cubeEdges}{%
  \draw[circuit] (A)--(B)--(F)--(E)--cycle;
  \draw[circuit] (B)--(C)--(G)--(F) (E)--(H)--(G);
  \draw[hiddenwire] (A)--(D)--(C) (D)--(H);}

\newcommand{\cubeProblemFigure}{%
\begin{center}
\begin{tikzpicture}[x=1cm,y=1cm]
  \cubeCoordinates\cubeEdges
  \foreach \P in {A,B,C,D,E,F,G,H}{\node[bluenode] at (\P) {};}
  \node[xeslabel,below left] at (A) {$A$};\node[xeslabel,below right] at (B) {$B$};
  \node[xeslabel,right] at (C) {$C$};\node[xeslabel,above right] at (G) {$C_1$};
  \node[xessmall] at (.30,-1.62) {虚线棱在立方体后方，所有投影交叉处均不另设节点};
\end{tikzpicture}
\end{center}}

\newcommand{\cubePotentialFigure}{%
\begin{center}
\begin{tikzpicture}[x=.69cm,y=.69cm]
  \begin{scope}[shift={(-6.0,0)}]
    \cubeCoordinates\cubeEdges
    \potentialNode{A}{1}{Gold!72}\potentialNode{B}{0}{Ink!20}
    \potentialNode{E}{\frac9{14}}{DeepBlue!28}\potentialNode{D}{\frac9{14}}{DeepBlue!28}
    \potentialNode{H}{\frac47}{DeepBlue!23}\potentialNode{F}{\frac5{14}}{DeepBlue!14}
    \potentialNode{C}{\frac5{14}}{DeepBlue!14}\potentialNode{G}{\frac37}{DeepBlue!18}
    \node[xestiny] at (.30,-1.70) {(1) 邻点 $A\to B$};
  \end{scope}
  \begin{scope}[shift={(0,0)}]
    \cubeCoordinates\cubeEdges
    \potentialNode{A}{1}{Gold!72}\potentialNode{C}{0}{Ink!20}
    \potentialNode{E}{\frac23}{DeepBlue!30}\potentialNode{D}{\frac12}{RuleGray!55}
    \potentialNode{H}{\frac12}{RuleGray!55}\potentialNode{B}{\frac12}{RuleGray!55}
    \potentialNode{F}{\frac12}{RuleGray!55}\potentialNode{G}{\frac13}{DeepBlue!13}
    \node[xestiny] at (.30,-1.70) {(2) 面对角 $A\to C$};
  \end{scope}
  \begin{scope}[shift={(6.0,0)}]
    \cubeCoordinates\cubeEdges
    \potentialNode{A}{1}{Gold!72}\potentialNode{G}{0}{Ink!20}
    \potentialNode{B}{\frac35}{DeepBlue!27}\potentialNode{D}{\frac35}{DeepBlue!27}
    \potentialNode{E}{\frac35}{DeepBlue!27}\potentialNode{C}{\frac25}{DeepBlue!15}
    \potentialNode{F}{\frac25}{DeepBlue!15}\potentialNode{H}{\frac25}{DeepBlue!15}
    \node[xestiny] at (.30,-1.70) {(3) 体对角 $A\to G$};
  \end{scope}
\end{tikzpicture}
\end{center}}

\newcommand{\infiniteLadderProblemFigure}{%
\begin{center}
\begin{tikzpicture}[x=1cm,y=1cm]
  \coordinate (a) at (-3.7,.75);\coordinate (b) at (-3.7,-.75);
  \node[port] at (a) {};\node[port] at (b) {};
  \node[xeslabel,left] at (a) {$a$};\node[xeslabel,left] at (b) {$b$};
  \draw[circuit] (a)--(-3.45,.75) (b)--(-3.45,-.75);
  \foreach \i in {0,1,2,3}{
    \pgfmathsetmacro{\xa}{-3.45+1.65*\i}\pgfmathsetmacro{\xb}{-2.05+1.65*\i}
    \zigres{\xa,.75}{\xb,.75}\zigres{\xa,-.75}{\xb,-.75}
    \zigres{\xb,.75}{\xb,-.75}
    \ifnum\i<3\draw[circuit] (\xb,.75)--({\xb+.25},.75) (\xb,-.75)--({\xb+.25},-.75);\fi}
  \node[xestiny,above] at (-2.75,1.02) {$R$};
  \draw[circuit,densely dotted] (2.90,.75)--(4.10,.75) (2.90,-.75)--(4.10,-.75);
  \node[xeslabel,right] at (4.10,0) {$\cdots$};
\end{tikzpicture}
\end{center}}

\newcommand{\ladderRecursionFigure}{%
\begin{center}
\begin{tikzpicture}[x=1cm,y=1cm]
  \coordinate (a) at (-3.4,.85);\coordinate (b) at (-3.4,-.85);
  \node[port] at (a) {};\node[port] at (b) {};\node[xeslabel,left] at (a) {$a$};\node[xeslabel,left] at (b) {$b$};
  \draw[circuit] (a)--(-3.15,.85) (b)--(-3.15,-.85);
  \zigres{-3.15,.85}{-1.65,.85}\zigres{-3.15,-.85}{-1.65,-.85}
  \node[xestiny,above] at (-2.40,1.12) {$R$};\node[xestiny,below] at (-2.40,-1.12) {$R$};
  \coordinate (u) at (-1.45,.85);\coordinate (d) at (-1.45,-.85);
  \draw[circuit] (-1.65,.85)--(u) (-1.65,-.85)--(d);
  \zigres{u}{d}\node[xestiny,left] at (-1.73,0) {$R$};
  \draw[circuit] (u)--(-.65,.85) (d)--(-.65,-.85);
  \draw[DeepBlue!70,dashed,rounded corners=3pt,line width=.8pt] (-.65,1.25) rectangle (2.65,-1.25);
  \node[xeslabel] at (1.0,0) {无限尾部 $\equiv X$};
  \draw[circuit] (2.65,.85)--(3.15,.85) (2.65,-.85)--(3.15,-.85);
  \node[port] at (3.15,.85) {};\node[port] at (3.15,-.85) {};
  \node[xessmall] at (0,-1.72) {$X=R+\bigl(R\parallel X\bigr)+R$};
\end{tikzpicture}
\end{center}}

\newcommand{\threeRailProblemFigure}{%
\begin{center}
\begin{tikzpicture}[x=1cm,y=1cm]
  \foreach \r/\lab in {1.35/A,0/B,-1.35/C}{
    \node[port] at (-4.05,\r) {};\node[xeslabel,left] at (-4.05,\r) {$\lab$};
    \draw[circuit] (-4.05,\r)--(-3.82,\r);
    \foreach \i in {0,1,2,3}{
      \pgfmathsetmacro{\xa}{-3.82+1.65*\i}\pgfmathsetmacro{\xb}{-2.48+1.65*\i}
      \zigres{\xa,\r}{\xb,\r}
      \ifnum\i<3\draw[circuit] (\xb,\r)--({\xb+.31},\r);\fi}}
  \foreach \i in {0,1,2,3}{
    \pgfmathsetmacro{\x}{-2.48+1.65*\i}
    \zigres{\x,1.35}{\x,0}\zigres{\x,0}{\x,-1.35}}
  \draw[circuit,densely dotted] (2.47,1.35)--(3.55,1.35) (2.47,0)--(3.55,0) (2.47,-1.35)--(3.55,-1.35);
  \node[xeslabel,right] at (3.55,0) {$\cdots$};
  \node[xestiny,above] at (-3.15,1.65) {$R$};\node[xestiny,right] at (-2.20,.68) {$R$};
\end{tikzpicture}
\end{center}}

\newcommand{\threeRailModesFigure}{%
\begin{center}
\begin{tikzpicture}[x=1cm,y=1cm]
  \begin{scope}[shift={(-3.1,0)}]
    \foreach \y in {1,0,-1}{\node[bluenode] at (0,\y) {};}
    \draw[circuit] (0,1)--(0,0)--(0,-1);
    \draw[currentarrow] (.35,0)--(.35,1.0);\draw[currentarrow] (.35,0)--(.35,-1.0);
    \node[xeslabel,right] at (.58,1) {$+1$};\node[xeslabel,right] at (.58,0) {$0$};\node[xeslabel,right] at (.58,-1) {$-1$};
    \node[xessmall] at (0,-1.52) {$\boldsymbol e_1=(1,0,-1)$，$\mu=1$};
  \end{scope}
  \begin{scope}[shift={(3.1,0)}]
    \foreach \y in {1,0,-1}{\node[bluenode] at (0,\y) {};}
    \draw[circuit] (0,1)--(0,0)--(0,-1);
    \draw[currentarrow] (.35,0)--(.35,.78);\draw[currentarrow] (.35,.0)--(.35,-.78);
    \draw[currentarrow] (-.35,.85)--(-.35,-.85);
    \node[xeslabel,right] at (.58,1) {$+1$};\node[xeslabel,right] at (.58,0) {$-2$};\node[xeslabel,right] at (.58,-1) {$+1$};
    \node[xessmall] at (0,-1.52) {$\boldsymbol e_3=(1,-2,1)$，$\mu=3$};
  \end{scope}
\end{tikzpicture}
\end{center}}

\newcommand{\nestedTriangleProblemFigure}{%
\begin{center}
\begin{tikzpicture}[x=1cm,y=1cm]
  \coordinate (O) at (0,0);
  \foreach \k in {0,...,6}{
    \pgfmathsetmacro{\rad}{2.55/(2^\k)}
    \pgfmathsetmacro{\ang}{90+180*mod(\k,2)}
    \draw[circuit,line width=.72pt]
      ($(O)+(\ang:\rad)$)--($(O)+({\ang+120}:\rad)$)--($(O)+({\ang+240}:\rad)$)--cycle;}
  \node[xeslabel,below left] at ($(O)+(210:2.55)$) {$A$};
  \node[xeslabel,below right] at ($(O)+(330:2.55)$) {$B$};
  \node[xeslabel,above] at ($(O)+(90:2.55)$) {$C$};
\end{tikzpicture}
\end{center}}

\newcommand{\nestedTriangleRecursionFigure}{%
\begin{center}
\begin{tikzpicture}[x=1cm,y=1cm]
  \coordinate (A) at (-2.55,-1.45);\coordinate (B) at (2.55,-1.45);\coordinate (C) at (0,2.97);
  \coordinate (Mab) at (0,-1.45);\coordinate (Mbc) at (1.275,.76);\coordinate (Mca) at (-1.275,.76);
  \draw[solutionwire] (A)--(Mab)--(B) (B)--(Mbc)--(C) (C)--(Mca)--(A);
  \draw[activewire] (Mab)--(Mbc)--(Mca)--cycle;
  \foreach \P in {A,B,C}{\node[goldnode] at (\P) {};}
  \foreach \P in {Mab,Mbc,Mca}{\node[bluenode] at (\P) {};}
  \node[xeslabel,below left] at (A) {$A$};\node[xeslabel,below right] at (B) {$B$};\node[xeslabel,above] at (C) {$C$};
  \node[xestiny,below] at (-1.28,-1.55) {$r/2$};\node[xestiny,below] at (1.28,-1.55) {$r/2$};
  \node[xestiny] at (0,.06) {内网等效支路 $kr/2$};
  \node[xessmall] at (0,-2.05) {蓝色三角形不是新增金属丝，而是对全部后续嵌套网络的三端等效};
\end{tikzpicture}
\end{center}}

\newcommand{\sierpinskiProblemFigure}{%
\begin{center}
\begin{tikzpicture}[x=1cm,y=1cm]
  \begin{scope}[shift={(-5.2,0)}]
    \fill[DeepBlue!48] (-1.45,-1.0)--(1.45,-1.0)--(0,1.51)--cycle;
    \draw[circuit] (-1.45,-1.0)--(1.45,-1.0)--(0,1.51)--cycle;
    \node[xessmall] at (0,-1.48) {初始};
  \end{scope}
  \begin{scope}[shift={(0,0)}]
    \foreach \x/\y in {-0.725/-1.0,0.725/-1.0,0/.255}{
      \fill[DeepBlue!48] (\x-.725,\y)--(\x+.725,\y)--(\x,{\y+1.255})--cycle;
      \draw[circuit] (\x-.725,\y)--(\x+.725,\y)--(\x,{\y+1.255})--cycle;}
    \node[xessmall] at (0,-1.48) {第一阶};
  \end{scope}
  \begin{scope}[shift={(5.2,0)}]
    \foreach \X/\Y in {-1.0875/-1.0,-.3625/-1.0,.3625/-1.0,1.0875/-1.0,-.725/-.3725,.725/-.3725,-.3625/.255,.3625/.255,0/.8825}{
      \fill[DeepBlue!48] (\X-.3625,\Y)--(\X+.3625,\Y)--(\X,{\Y+.6275})--cycle;
      \draw[circuit,line width=.6pt] (\X-.3625,\Y)--(\X+.3625,\Y)--(\X,{\Y+.6275})--cycle;}
    \node[xessmall] at (0,-1.48) {第二阶};
  \end{scope}
  \node[xestiny] at (-6.75,-1.10) {$A$};\node[xestiny] at (-3.65,-1.10) {$B$};\node[xestiny] at (-5.20,1.68) {$C$};
\end{tikzpicture}
\end{center}}

\newcommand{\sierpinskiResistanceFigure}{%
\begin{center}
\begin{tikzpicture}[x=1cm,y=1cm]
  \coordinate (A) at (-3.0,-1.25);\coordinate (B) at (3.0,-1.25);\coordinate (C) at (0,3.95);
  \coordinate (Mab) at (0,-1.25);\coordinate (Mbc) at (1.50,1.35);\coordinate (Mca) at (-1.50,1.35);
  \foreach \P/\Q/\R in {A/Mab/Mca,Mab/B/Mbc,Mca/Mbc/C}{
    \draw[solutionwire,fill=DeepBlue!5] (\P)--(\Q)--(\R)--cycle;}
  \foreach \P in {A,B,C,Mab,Mbc,Mca}{\node[bluenode] at (\P) {};}
  \node[xeslabel,below left] at (A) {$A$};\node[xeslabel,below right] at (B) {$B$};\node[xeslabel,above] at (C) {$C$};
  \node[xessmall] at (-1.5,-1.95) {缩小的同型三端网};
  \draw[->,TextGray,line width=.6pt] (-1.5,-1.70)--(-1.5,-.55);
  \node[eqbox,anchor=west] at (3.65,1.35) {$\displaystyle
    R_{n+1}=\frac53R_n\quad(\text{单段电阻不缩放})$};
  \node[eqbox,anchor=west] at (3.65,.20) {$\displaystyle
    r_{n+1}=\frac12r_n\quad(\text{真实长度缩放})$};
  \node[eqbox,anchor=west] at (3.65,-.95) {$\displaystyle
    R_{n+1}^{\rm phys}=\frac56R_n^{\rm phys}$};
\end{tikzpicture}
\end{center}}


原则上，基尔霍夫定律足以处理绝大多数集中参数电路；但网络一旦变复杂，直接设置每条支路电流、
列出全部节点和回路方程，会迅速产生规模很大的线性方程组。这样的做法可靠，却冗长、难以检查，也很
容易在符号上出错。因此，对只含线性元件的网络，应优先利用线性叠加、等效电源、对称性、电桥平衡、
星形--三角形变换和自相似性，把网络先化简，再在必要时使用基尔霍夫方程。对于计算机，求解大规模
线性方程组并不困难；对人工竞赛解题而言，找到结构才是关键。


\subsection{叠加原理}


节点电压法、支路电流法和基尔霍夫方程不仅适用于电阻，也可与任意元件的伏安关系联立。若元件是
非线性的，得到的通常是非线性方程组，叠加原理一般不成立；若网络只含线性电阻和独立电源，方程组对
电源产生的非齐次项具有可加性。于是多个独立电源同时工作时，任一支路电流或任一节点电势，等于各个
电源单独作用时对应结果的代数和。

考察某一个独立电源的贡献时，其余独立电源要“置零”，同时保留它们的内阻和原有连接关系：

\begin{itemize}
  \item 理想电压源置零，即把其电压改为零，等效为短路；
  \item 理想电流源置零，即把其电流改为零，等效为断路。
\end{itemize}

下图把本章前面含电压源 $E_0$、电流源 $I_0$ 和两个电阻的电路分成两个子问题。仅保留电压源时，
$R_1,R_2$ 串联，流过 $R_2$ 的电流为
\[
  I_1=\frac{E_0}{R_1+R_2}.
\]
仅保留电流源时，$R_1,R_2$ 并联，按分流关系，流过 $R_2$ 的电流为
\[
  I_2=\frac{R_1}{R_1+R_2}I_0.
\]
若图示两部分在 $R_2$ 中产生的电流方向相同，则原电路中
$I_{R_2}=I_1+I_2$；若方向相反，就应按所选正方向作代数相加。

\par\noindent\begin{minipage}{\linewidth}
\superpositionFigure
\diagramcaption{叠加原理。左图只保留电压源，右图只保留电流源；被置零的电压源短路，被置零的电流源断路。}\label{ch21:fig:s37-6-1}
\end{minipage}\par\medskip

叠加的是电压和电流，不是功率。由于 $P=I^2R$ 或 $P=U^2/R$ 是二次量，通常不能把各电源单独作用时
的功率直接相加。


\subsection{等效电源定理}


线性叠加还有两个重要推论：戴维南（Thevenin）等效电压源定理和诺顿（Norton）等效电流源定理。
无论内部结构多复杂，一个线性有源二端网络对外只表现为端电压与端电流之间的一条直线关系；只要确定
这条直线的截距和斜率，就不必再追踪网络内部的每条支路。

对图~\ref{ch21:fig:s37-6-2}(a) 的二端网络，规定测试电流 $i_t$ 流入上端口，端电压为 $v_t$。在端口接入可调理想电流源，
得到图~\ref{ch21:fig:s37-6-2}(b)。由叠加原理，$v_t$ 可拆成两部分：

\begin{enumerate}
  \item 保留测试电流源，把网络内部所有独立源置零。此时内部成为纯电阻网络，端口等效电阻记为 $r$，
  因而图~\ref{ch21:fig:s37-6-2}(c) 中 $v_a=i_t r$；
  \item 断开测试电流源，恢复网络内部全部独立源。端口电流为零，图~\ref{ch21:fig:s37-6-2}(d) 中的端电压 $v_b$ 就是
  开路电压。
\end{enumerate}

所以
\[
  v_t=v_a+v_b=i_t r+v_b.
\]
这正是一个理想电压源与内阻串联时的端口关系。

\par\noindent\begin{minipage}{\linewidth}
\theveninDerivationFigure
\diagramcaption{戴维南定理的推导。四幅图的端口方向完全一致，第三幅只保留测试源，第四幅只保留网络内部独立源。}\label{ch21:fig:s37-6-2}
\end{minipage}\par\medskip

具体地说，任意线性有源二端网络都可等效为开路电压 $V_{\rm Th}$ 与电阻 $R_{\rm Th}$ 串联：
$V_{\rm Th}$ 由原网络端口开路时的电压得到；$R_{\rm Th}$ 由内部所有独立源置零后、从端口向里看见的
等效电阻得到。若网络中还含受控源，受控源不能置零，应在端口加测试源求 $R_{\rm Th}=v_t/i_t$。

同一网络也可等效为短路电流 $I_{\rm N}$ 与电阻 $R_{\rm N}$ 并联。$I_{\rm N}$ 是端口短路时的电流，
$R_{\rm N}$ 与 $R_{\rm Th}$ 相同。两种等效电路在端口上满足
\begin{equation}
  \boxed{R_{\rm Th}=R_{\rm N},\qquad V_{\rm Th}=I_{\rm N}R_{\rm N}.}
\end{equation}

\par\noindent\begin{minipage}{\linewidth}
\sourceEquivalenceFigure
\diagramcaption{同一个线性有源二端网络的戴维南形式与诺顿形式。等效只保证端口外部特性相同。}\label{ch21:fig:s37-6-3}
\end{minipage}\par\medskip

“等效”并不意味着三种电路内部的电流和功率分布相同；它只表示接到同一个负载后，端口电压和负载电流
完全一致。


% EM teaching enrichment 20260927
\paragraph{负载功率为何在阻值相等时最大}
当等效内阻 $R_{\rm Th}>0$、外接纯电阻 $R_L>0$ 时，
\[
 P_L=\frac{V_{\rm Th}^2R_L}{(R_{\rm Th}+R_L)^2},\qquad
 \frac{\dd P_L}{\dd R_L}
 =\frac{V_{\rm Th}^2(R_{\rm Th}-R_L)}{(R_{\rm Th}+R_L)^3}.
\]
故 $R_L=R_{\rm Th}$ 时最大，$P_{L,\max}=V_{\rm Th}^2/(4R_{\rm Th})$。这里最大化的是负载功率，不是输电效率；在这个简单串联等效模型中，匹配时负载与内阻耗散相等，效率仅为一半。端口等效也不能据此还原复杂原网络内部每个元件的耗散。


\subsection{对称性}


若交换网络中的两个节点后，所有连接、元件参数以及外加电源条件都保持不变，那么这两个节点不可区分，
电势必然相等。等势节点之间没有电流，可以按需要用理想导线相连，也可以把连接它们的支路删去。使用
对称性时，必须把“电路本身的对称”与“端口激励的对称”同时检查；只看几何外形相似而忽略电流从哪里
流入、从哪里流出，是复杂电阻网络中最常见的错误之一。


% EM teaching enrichment 20260927
\paragraph{等势结论来自同一个边值问题}
把端口 $A$ 固定为 $U$、$B$ 固定为零。若一个变换保持两个端口各自不动，并把等阻支路一一对应，变换前后的节点电势都满足同一组 KCL；由电阻网络解的唯一性，相互交换的节点必等势。
若变换交换 $A,B$，一般得到的是 $V_{i'}=U-V_i$，而非 $V_{i'}=V_i$。只有处在该变换不动位置的节点才能推出 $V=U/2$。
例如所有棱阻值为 $R$ 的正方体，从一对体对角顶点通电，靠近入端的三个顶点等势，靠近出端的三个也等势，三层等效为
\[
 R_{AB}=\frac R3+\frac R6+\frac R3=\frac{5R}{6}.
\]
若把出端换到相邻顶点，原先的三点分组不再保持端口条件，不能照搬上述结果。


\subsection{电桥平衡}


图~\ref{ch21:fig:s37-6-4} 中，电源跨接在左右两端。若桥臂满足
\[
  \boxed{R_1R_4=R_2R_3},
\]
则上、下两个中点电势相等，$R_5$ 中电流为零，这就是电桥平衡。此时 $R_5$ 既可看作没有电流的断路支路
而删去，也可利用两端等势把它短接；两种处理对其余电路的电流、电压都没有影响。要注意，“断路”说的是
$R_5$ 中无电流，“短路”说的是它的两个端点等势，二者描述的是同一平衡状态的不同侧面。

\par\noindent\begin{minipage}{\linewidth}
\bridgeFigure
\diagramcaption{平衡电桥。$R_5$ 的上下端分别是两个分压支路的中点；平衡条件下二者等势。}\label{ch21:fig:s37-6-4}
\end{minipage}\par\medskip


% EM teaching enrichment 20260927
\paragraph{先断开桥支路，比较两个分压}
设图中上支路由左至右为 $R_1,R_2$，下支路为 $R_3,R_4$，左右端电势分别为 $U,0$。暂断开 $R_5$，两中点电势为
\[
 V_{\rm up}=U\frac{R_2}{R_1+R_2},\qquad
 V_{\rm down}=U\frac{R_4}{R_3+R_4}.
\]
相等即给出 $R_1R_4=R_2R_3$。接回任意有限 $R_5$ 不产生桥电流，因而原分压保持成立；若不相等，接回后两个分压都会受加载影响。
还可把桥端口外的网络化成戴维南源，其内阻为
$(R_1\parallel R_2)+(R_3\parallel R_4)$，桥电流便等于开路中点电压差除以该内阻与 $R_5$ 之和。这给出了偏离平衡时可继续计算的路线。


% EM structure and teaching revision 20260928
\subsection{端口等效的测定与适用边界}
线性有源二端网络可在端口写成 $U=U_{\mathrm{oc}}-IR_{\mathrm{th}}$，其中 $I$ 取网络向负载输出的方向。开路电压确定截距，短路电流在存在且有限时给出 $R_{\mathrm{th}}=U_{\mathrm{oc}}/I_{\mathrm{sc}}$。
只有独立源的电路，求等效电阻时才可把理想独立电压源短接、理想独立电流源断开；实际源内阻仍应保留。含受控源时必须保留控制关系，并在端口加测试电压或测试电流来求比值。
对 $R_{\mathrm{th}}>0$ 的被动负载 $R_L$，
\[
 P_L=\frac{U_{\mathrm{oc}}^2R_L}{(R_{\mathrm{th}}+R_L)^2},\qquad
 \frac{\dd P_L}{\dd R_L}=\frac{U_{\mathrm{oc}}^2(R_{\mathrm{th}}-R_L)}{(R_{\mathrm{th}}+R_L)^3}.
\]
最大功率对应阻值匹配，但此时等效内阻与负载各消耗一半功率。若追求输电效率，应使线路电阻远小于负载的等效阻值，目标与最大功率传输不同。
网络对称性也必须连同端口激励一起判断：图形上对称的两个节点，在不对称接入电源后未必等势。


\subsection{专题例题与方法应用}

\begin{origexamplebox}{【电流源】}\label{ch21:ex:s37-6-1}
小明利用一个电流源向外供电，但是小红却误以为是一个电压源。小明发现无法说服小红，即电流源和
电压源是等价的。试帮助小明证明这个结论，并利用这个关系求解电压源并联的公式。
\end{origexamplebox}\par\nopagebreak[4]

\par\noindent\begin{minipage}{\linewidth}
\sourceTransformExampleFigure
\diagramcaption{解图：电压源串联内阻与电流源并联内阻的端口方向。两图对同一负载给出相同的 $U,I$。}
\end{minipage}\par\medskip

\begin{xesolution}
这里的“等价”指\textbf{带内阻电源的端口等价}，并不是说孤立的理想电压源与孤立的理想电流源可以
无条件互换。设实际电压源的电动势为 $E$、串联内阻为 $r$，负载端电压为 $U$，向负载输出的电流为
$I$。由 KVL，端口关系为
\[
  U=E-Ir,
  \qquad I=\frac Er-\frac Ur.
\]
若改用输出电流 $I_{\rm N}=E/r$ 的理想电流源，并在端口并联同一个电阻 $r$，KCL 给出
\[
  I=I_{\rm N}-\frac Ur=\frac Er-\frac Ur,
\]
与上式完全相同。因此
\[
  \boxed{E\ \text{串联}\ r
  \quad\Longleftrightarrow\quad
  I_{\rm N}=\frac Er\ \text{并联}\ r}.
\]

现在把 $n$ 个电压源 $(E_i,r_i)$ 同极性并联。先把每一支路变成诺顿形式：电流源 $E_i/r_i$ 与
$r_i$ 并联。并联电流源的电流相加，并联内阻的电导相加，故
\[
  I_{\rm N,eq}=\sum_{i=1}^n\frac{E_i}{r_i},
  \qquad
  \frac1{r_{\rm eq}}=\sum_{i=1}^n\frac1{r_i}.
\]
再变回电压源，得到
\[
  \boxed{r_{\rm eq}=\left(\sum_{i=1}^n\frac1{r_i}\right)^{-1}},
  \qquad
  \boxed{E_{\rm eq}=r_{\rm eq}\sum_{i=1}^n\frac{E_i}{r_i}}.
\]
对两个电源，公式化为
\[
  \boxed{r_{\rm eq}=\frac{r_1r_2}{r_1+r_2},\qquad
  E_{\rm eq}=\frac{E_1r_2+E_2r_1}{r_1+r_2}}.
\]
若某支路极性相反，应把相应 $E_i$ 取负。两个电动势不同、内阻又被理想化为零的电压源不能直接并联，
因为模型会要求无限大的环流；真实电源的内阻正是公式成立所必需的物理条件。
\end{xesolution}



\begin{origexamplebox}{【田电阻】}\label{ch21:ex:s37-6-2}
有一个田字格的电阻网络，每条边电阻为 $R$，试求：

(1) 对角线两个点 $A,C$ 的等效电阻；

(2) 相邻的两个顶点 $A,B$ 的等效电阻。
\par\noindent\begin{minipage}{\linewidth}
\gridProblemFigure
\diagramcaption{九个交点都是节点，十二个相邻节点间的线段各代表一个电阻 $R$。}\label{ch21:fig:s37-6-6}
\end{minipage}\par\medskip
\end{origexamplebox}\par\nopagebreak[4]


\begin{xesolution}
\textbf{(1) 对角端点 $A,C$。} 令 $\varphi_A=V,\ \varphi_C=0$。关于对角线 $AC$ 反射，成对节点
等势；再结合绕中心旋转 $180^\circ$ 后交换 $A,C$，可知中心点和另两个外角的电势均为 $V/2$。
与 $A$ 相邻的两个边中点等势，设其电势为 $x$。其中任一点连接 $A$ 和两个 $V/2$ 节点，KCL 给出
\[
  \frac{x-V}{R}+2\frac{x-V/2}{R}=0,
  \qquad x=\frac{2V}{3}.
\]
从 $A$ 流出的总电流是
\[
  I=2\frac{V-x}{R}=\frac{2V}{3R},
\]
故
\[
  \boxed{R_{AC}=\frac VI=\frac32R}.
\]

\par\noindent\begin{minipage}{\linewidth}
\gridPotentialDiagonalFigure
\diagramcaption{解图（1）：把 $A,C$ 的电势分别归一化为 $1,0$。同色节点等势，标注放在网格空隙内而不遮挡支路。}
\end{minipage}\par\medskip

\textbf{(2) 相邻外角 $A,B$。} 仍令 $\varphi_A=V,\ \varphi_B=0$。左右反射同时交换两个端点，
所以竖直中线上的三个节点均为 $V/2$。设左边中点电势为 $u$、左上角为 $w$，其右侧镜像节点的
电势分别为 $V-u,V-w$。左上角和左边中点的 KCL 为
\[
  2w-u-\frac V2=0,
  \qquad
  3u-V-w-\frac V2=0.
\]
解得 $u=7V/10,\ w=3V/5$。从 $A$ 流出的电流
\[
  I=\frac{V-V/2}{R}+\frac{V-7V/10}{R}=\frac{4V}{5R},
\]
因此
\[
  \boxed{R_{AB}=\frac54R}.
\]

\par\noindent\begin{minipage}{\linewidth}
\gridPotentialAdjacentFigure
\diagramcaption{解图（2）：相邻端点激励不再具有对角情形的全部对称性；图中的九个电势已经逐点满足 KCL。}
\end{minipage}\par\medskip
\end{xesolution}



\begin{origexamplebox}{【正四面体电阻】}\label{ch21:ex:s37-6-3}
有一个正四面体框架的电阻网络，每条边电阻为 $R$，试求任意两点间的等效电阻。
\par\noindent\begin{minipage}{\linewidth}
\tetrahedronProblemFigure
\diagramcaption{四个顶点、六条棱均完整画出；交叉处没有节点。}\label{ch21:fig:s37-6-7}
\end{minipage}\par\medskip
\end{origexamplebox}\par\nopagebreak[4]


\begin{xesolution}
从任意两顶点 $A,B$ 接入电源。其余两顶点 $C,D$ 在保持 $A,B$ 不动的对称变换下可以交换，所以
$\varphi_C=\varphi_D$，连接 $C,D$ 的电阻中没有电流，可删去。此后 $A$ 到 $B$ 之间有三条并联支路：
一条是直接的 $R$，另两条分别是 $A-C-B$ 和 $A-D-B$，每条为 $2R$。于是
\[
  \frac1{R_{AB}}=\frac1R+\frac1{2R}+\frac1{2R}=\frac2R,
  \qquad
  \boxed{R_{AB}=\frac R2}.
\]
这个结论与所选顶点无关，因为正四面体的任意两顶点在网络拓扑上完全等价。
\end{xesolution}



\begin{origexamplebox}{【$N$ 边形电阻】}\label{ch21:ex:s37-6-4}
有一个 $N$ 边形电阻网络（图中 $N=5$），每条边、每条对角线都用阻值为 $R$ 的电阻相连，试求
任意两点间的等效电阻。
\par\noindent\begin{minipage}{\linewidth}
\completeGraphProblemFigure
\diagramcaption{图示为 $K_5$；一般情形是任意两顶点之间都有且只有一个电阻 $R$。}\label{ch21:fig:s37-6-8}
\end{minipage}\par\medskip
\end{origexamplebox}\par\nopagebreak[4]


\begin{xesolution}
取所求端点为 $A,B$，加电压 $V$。其余 $N-2$ 个顶点在保持 $A,B$ 不变的任意置换下都等价，故电势
相同。它们之间虽有大量电阻，但每一条两端等势，电流均为零。由交换 $A,B$ 的反对称性，这些中间点
电势都是 $V/2$。

从 $A$ 流出的电流包括直接流向 $B$ 的 $V/R$，以及流向 $N-2$ 个中间点的
$(N-2)V/(2R)$：
\[
  I=\frac VR+(N-2)\frac{V}{2R}=\frac{NV}{2R}.
\]
所以
\[
  \boxed{R_{AB}=\frac{2R}{N}}.
\]
代入 $N=4$ 就得到上一题正四面体的 $R/2$；二者本质上分别是完全图 $K_N$ 与 $K_4$。

\par\noindent\begin{minipage}{\linewidth}
\completeGraphReductionFigure
\diagramcaption{解图：固定端点 $A,B$ 后，其余 $N-2$ 点全部等势。虚线框表示把这些点视为同一电势层，而非新增导线。}
\end{minipage}\par\medskip
\end{xesolution}



\begin{origexamplebox}{【变种八边形电阻】}\label{ch21:ex:s37-6-5}
有一个八边形电阻网络，除最长的对角线外，每条边、每条对角线都用阻值为 $R$ 的电阻相连，试求相距
最远的两点间的等效电阻。
\par\noindent\begin{minipage}{\linewidth}
\octagonProblemFigure
\diagramcaption{八边形中唯独四条直径不接电阻；任意其他两顶点之间均接一个 $R$。}\label{ch21:fig:s37-6-9}
\end{minipage}\par\medskip
\end{origexamplebox}\par\nopagebreak[4]


\begin{xesolution}
相距最远的端点 $A,B$ 互为对顶点，它们之间恰好没有直径支路。设端电势为 $V,0$，试令其余六点均为
$V/2$：任一中间点与 $A,B$ 相连的两支路电流相互抵消，与其他中间点相连的支路电流为零，所以各点
KCL均满足。正电阻网络在端电势给定后内部电势解唯一，故这就是实际电势，不需要声称任意交换顶点都
保持网络不变。每个中间点都通过一个 $R$ 接到 $A$，并通过另一个 $R$ 接到
$B$；中间点之间的支路两端等势，均无电流。

网络于是化成六条相同的 $A$--中间点--$B$ 支路并联，每条阻值 $2R$：
\[
  \boxed{R_{AB}=\frac{2R}{6}=\frac R3}.
\]
这里不能误加一条 $A,B$ 间的直接电阻；“最长对角线除外”正是本题与完全图 $K_8$ 的区别。

\par\noindent\begin{minipage}{\linewidth}
\octagonReductionFigure
\diagramcaption{解图：六个蓝色中间点等势，内部连线均不载流；真正参与端口电流的是六条长度为两个电阻的并联通路。}
\end{minipage}\par\medskip
\end{xesolution}



\begin{origexamplebox}{【正方体电阻】}\label{ch21:ex:s37-6-6}
有一个立方体电阻网络，每条棱电阻为 $R$。求：

(1) 相邻两顶点；

(2) 面对角线两顶点；

(3) 体对角线两顶点的等效电阻。
\par\noindent\begin{minipage}{\linewidth}
\cubeProblemFigure
\diagramcaption{实线和虚线共同构成立方体的十二条棱；虚线只表示被遮挡的三条棱。}\label{ch21:fig:s37-6-10}
\end{minipage}\par\medskip
\end{origexamplebox}\par\nopagebreak[4]


\begin{xesolution}
统一令输入端电势为 $V$、输出端为 $0$。下图列出三种接法中所有顶点的归一化电势，数值由对称性合并
节点后写 KCL 得到；把每个内部节点的相邻电势代入，都可直接验证流入等于流出。

\par\noindent\begin{minipage}{\linewidth}
\cubePotentialFigure
\diagramcaption{解图：三种端点选择的归一化电势。每个数值紧邻顶点外侧，虚线棱与实线棱的空间关系保持一致。}
\end{minipage}\par\medskip

\textbf{(1) 相邻顶点。} 用三位坐标 $000$ 至 $111$ 标记立方体顶点，一条棱只改变其中一位。
取 $000$ 为输入端、$100$ 为输出端。交换后两坐标可知 $010,001$ 等势，记为 $x$；
反转第一坐标同时交换端电势，可知 $110,101$ 电势为 $V-x$。另记 $011$ 电势为 $y$，则 $111$ 为 $V-y$。
在 $010$ 和 $011$ 两点分别列KCL：
\[
 (x-V)+(x-(V-x))+(x-y)=0,\qquad
 2(y-x)+(y-(V-y))=0.
\]
整理为 $4x-y=2V,\ 4y-2x=V$，消去 $y$ 得 $14x=9V$，再得 $y=4V/7$。
因此输入端除直达输出端的一条棱外，另两条棱通向电势均为 $9V/14$ 的节点，
所以
\[
  I=\frac VR+2\frac{V-9V/14}{R}=\frac{12V}{7R},
  \qquad
  \boxed{R_{\text{邻}}=\frac{7R}{12}}.
\]

\textbf{(2) 面对角线顶点。} 取 $000$ 输入、$110$ 输出。交换前两坐标以及同时反转这两位，
可知 $100,010,101,011$ 电势均为 $V/2$。令 $001$ 的电势为 $x$，它的邻点为 $000,101,011$，故
\[
 (x-V)+2(x-V/2)=0,\qquad x=2V/3.
\]
输入端的两个同面邻点电势均为 $V/2$，第三个邻点电势为 $2V/3$，故
\[
  I=2\frac{V-V/2}{R}+\frac{V-2V/3}{R}=\frac{4V}{3R},
  \qquad
  \boxed{R_{\text{面}}=\frac{3R}{4}}.
\]

\textbf{(3) 体对角线顶点。} 输入端的三个邻点等势，设为 $x$；输出端的三个邻点由反对称性为 $V-x$。
任取一个 $x$ 节点，它连接输入端和两个 $V-x$ 节点，故
\[
  \frac{x-V}{R}+2\frac{x-(V-x)}R=0,
  \qquad x=\frac{3V}{5}.
\]
于是
\[
  I=3\frac{V-3V/5}{R}=\frac{6V}{5R},
  \qquad
  \boxed{R_{\text{体}}=\frac{5R}{6}}.
\]
三种结果依次增大，但这不是由“几何距离”直接决定的，而是由端点选择后整个网络允许的并联通道决定。
\end{xesolution}



\section{多端、无限与自相似网络}

\subsection{电流分布法}


求两点间等效电阻时，常规做法是在两点间加电压 $U$，求出总电流 $I$，再用 $R_{\rm eq}=U/I$。电流
分布法则反过来：从一点灌入给定电流 $I$，从另一点抽出同样的电流，先利用对称性和 KCL 确定各支路
如何分流，再求两点间电势差 $U$，最后用 $R_{\rm eq}=U/I$。

这种方法同样服从叠加。例如，要得到“电流从 $A$ 进、从 $B$ 出”的分布，可以把“从 $A$ 进、从
$C$ 出”和“从 $C$ 进、从 $B$ 出”两种分布按相同电流强度相加。中间节点 $C$ 的注入与抽出相互抵消，
剩下的正是所求端口条件。


% EM teaching enrichment 20260927
\paragraph{分流必须同时满足回路电压条件}
只写 KCL 可以构造许多看似守恒的支路电流，却未必对应真实电势。得到候选分流后，还应检查任意独立回路的 $\sum IR=0$。
一个实用核验是功率：若从 $A$ 注入 $I$、从 $B$ 抽出 $I$，则
\[
 I(V_A-V_B)=\sum_{\text{各电阻支路}}R_kI_k^2,
 \qquad R_{AB}=\frac{\sum_kR_kI_k^2}{I^2}.
\]
每条支路只计一次。对三角形三个等阻 $R$、端口取任意两顶点，直连支路分得 $2I/3$，绕行两阻各为 $I/3$，代入得 $R_{AB}=2R/3$，与 $R\parallel2R$ 相符。


\subsection{多端纯电阻网络}


只有 $N$ 个端口能与外界连接、内部只含电阻的电路，称为 $N$ 端纯电阻网络。一般的无源线性 $N$ 端
电阻网络，可以用每两端口之间各接一个适当电阻的完全图表示。端口若具有对称性，若干等效电阻可能相等；
若无对称性，完全图中的各支路阻值一般不同。

三端网络是最重要的特例。它既可写成三角形（$\Delta$）连接，也可写成星形（$Y$）连接。设端口电势
为 $\varphi_a,\varphi_b,\varphi_c$，星形中心电势为 $\varphi_0$，三条星形支路为 $R_a,R_b,R_c$。
中心节点满足 KCL：
\begin{equation}
  \frac{\varphi_a-\varphi_0}{R_a}
  +\frac{\varphi_b-\varphi_0}{R_b}
  +\frac{\varphi_c-\varphi_0}{R_c}=0.
\end{equation}
解得
\begin{equation}
  \varphi_0=
  \frac{\dfrac{\varphi_a}{R_a}+\dfrac{\varphi_b}{R_b}+\dfrac{\varphi_c}{R_c}}
       {\dfrac1{R_a}+\dfrac1{R_b}+\dfrac1{R_c}}
  =\frac{\varphi_aR_bR_c+\varphi_bR_cR_a+\varphi_cR_aR_b}
  {R_aR_b+R_bR_c+R_cR_a}.
\end{equation}
由端口 $A$ 流入星形网络的电流为
\begin{equation}
  I_a=\frac{\varphi_a-\varphi_0}{R_a}
  =\frac{\varphi_a(R_b+R_c)-\varphi_bR_c-\varphi_cR_b}
  {R_aR_b+R_bR_c+R_cR_a}
  =\frac{\varphi_a-\varphi_b}{R_{AB}}
  +\frac{\varphi_a-\varphi_c}{R_{AC}}.
\end{equation}
最后一个等号是三角形连接的端口电流。令
\[
  S=R_aR_b+R_bR_c+R_cR_a,
\]
比较 $\varphi_a,\varphi_b,\varphi_c$ 的系数，并循环置换端口，可得星形变三角形：
\[
  \boxed{R_{AB}=\frac{S}{R_c},\qquad
  R_{AC}=\frac{S}{R_b},\qquad
  R_{BC}=\frac{S}{R_a}}.
\]
反解得到三角形变星形：
\[
\boxed{
  R_a=\frac{R_{AB}R_{AC}}{R_{AB}+R_{AC}+R_{BC}},\quad
  R_b=\frac{R_{AB}R_{BC}}{R_{AB}+R_{AC}+R_{BC}},\quad
  R_c=\frac{R_{AC}R_{BC}}{R_{AB}+R_{AC}+R_{BC}}.}
\]

\par\noindent\begin{minipage}{\linewidth}
\deltaStarFigure
\diagramcaption{同一三端网络的三角形连接与星形连接。对应关系由端口字母决定，不应靠电阻在图中的朝向记忆。}\label{ch21:fig:s37-6-5}
\end{minipage}\par\medskip

星形--三角形变换能把许多表面上既不串联也不并联的网络改写成可逐级化简的形式；但在高度对称或
自相似网络中，先用等势、分流和递归往往更简洁。

% EM teaching enrichment 20260927
\paragraph{用三个端口状态检验星三角变换}
若三角形三边均为 $R$，对应星形每臂为 $R/3$。只在 $A,B$ 之间通电、$C$ 外端开路时，星形等效为 $2R/3$，恰与三角形的 $R\parallel2R$ 相同。
但若把 $C$ 接地，它就不再是开路条件，不能把通向 $C$ 的星形臂删掉。三端等效保证任意端电势下的三端电流均相同，比只匹配某一对端口的一个电阻值更强。
复杂网络宜先识别可局部替换的三端块，替换时保留端口名字；随后重新判断串、并联关系，避免按图形朝向误配公式分母。




% EM structure and teaching revision 20260928
\subsection{无限网络的递推、选根与收敛}
自相似结构允许用整体的等效量表示去掉首段后剩下的部分。设每段先串联一个 $R$，再把一个 $R$ 对地分流，余下网络等效仍为 $X$，则
\[
 X=R+\frac{RX}{R+X},\qquad X^2-RX-R^2=0.
\]
正电阻网络取 $X=(1+\sqrt5)R/2$。负根虽满足代数方程，却不满足被动端口的非负耗散条件。
更可靠的定义来自有限截断：从某个非负末端电阻 $X_0$ 出发，令 $X_{n+1}=R+RX_n/(R+X_n)$。一步后有 $R\le X_n\le2R$，且该区间内导数
$f'(X)=R^2/(R+X)^2\le1/4$，所以相邻近似间的差不断缩小，极限确为所选正根。
有多个端口时，一个标量等效电阻往往不够。应保留端口电势和注入电流的线性关系，或者利用对称性把端口激励分解成独立模式。未经证明就把三端网络替换成两端电阻，会丢失约束，使看似正确的递推求出错误结果。


\subsection{专题例题与方法应用}

\begin{origexamplebox}{【无穷网络】}\label{ch21:ex:s37-6-7}
求图示无限电阻网络中 $a,b$ 两点间的等效电阻。图中每个电阻的阻值均为 $R$。
\par\noindent\begin{minipage}{\linewidth}
\infiniteLadderProblemFigure
\diagramcaption{上下两条横向支路和每一根竖直梯档都为 $R$，网络向右无限延伸。}\label{ch21:fig:s37-6-11}
\end{minipage}\par\medskip
\end{origexamplebox}\par\nopagebreak[4]


\begin{xesolution}
设整个无限梯形网络的等效电阻为 $X$。删去最左边一格后，余下部分与原网络完全相同，因此从第一对
竖直节点向右看进去仍是 $X$。在这两个节点之间，第一根竖直电阻 $R$ 与余下网络 $X$ 并联；它们外侧
还各串联一个横向电阻 $R$。所以
\[
  X=2R+\frac{RX}{R+X}.
\]

\par\noindent\begin{minipage}{\linewidth}
\ladderRecursionFigure
\diagramcaption{解图：虚线框内的无限尾部仍等效为 $X$。两只横向 $R$ 串在 $R\parallel X$ 的两侧。}
\end{minipage}\par\medskip

令 $x=X/R$，有
\[
  x=2+\frac{x}{1+x},
  \qquad x^2-2x-2=0.
\]
电阻必须为正，舍去 $x=1-\sqrt3$，得到
\[
  \boxed{R_{ab}=X=(1+\sqrt3)R}.
\]
回代可见 $X>2R$，与“入口处必经上下两个 $R$，中间还有有限并联电阻”的结构相符。
\end{xesolution}



\begin{origexamplebox}{【三端无穷网络】}\label{ch21:ex:s37-6-8}
求图示无限电阻网络中 $A,C$ 两点和 $A,B$ 两点之间的等效电阻，已知图中每个电阻值均为 $R$。
\par\noindent\begin{minipage}{\linewidth}
\threeRailProblemFigure
\diagramcaption{题图示意：三条横向链向右无限延伸，相邻两行在每一列以竖直电阻连接；左端为 $A,B,C$。}
\end{minipage}\par\medskip
\end{origexamplebox}\par\nopagebreak[4]


\begin{xesolution}
这个网络有三个端口，不能只用一个标量电阻表示尾部。利用上下反射对称，把每一列的三个电势分解成
彼此独立的三种横向模式：
\[
  \boldsymbol e_0=(1,1,1),\qquad
  \boldsymbol e_1=(1,0,-1),\qquad
  \boldsymbol e_3=(1,-2,1).
\]
这些模式可直接从KCL导出。令一列的电势向量为 $\boldsymbol V=(V_A,V_B,V_C)^{\mathsf T}$，
流向本列两根竖直电阻的电流向量为
\[
 \boldsymbol I_{\rm vertical}=\frac1R
 \begin{pmatrix}1&-1&0\\-1&2&-1\\0&-1&1\end{pmatrix}\boldsymbol V
 =\frac1R L\boldsymbol V.
\]
逐项相乘得到 $L\boldsymbol e_0=0$、$L\boldsymbol e_1=\boldsymbol e_1$、
$L\boldsymbol e_3=3\boldsymbol e_3$。横向三根电阻同为 $R$，不会混合这些正交模式，
因此后两种模式分别看到无量纲横向电导 $\mu=1,3$。
端口一注入一抽出、总电流为零，不激发沿链流动的共同模式；各列相同的共同电势只表示参考电势的选择。

\par\noindent\begin{minipage}{\linewidth}
\threeRailModesFigure
\diagramcaption{解图：三行节点的两个有效电势模式。箭头长度表示相对电势（或注入电流）分量，不是实际支路。}
\end{minipage}\par\medskip

设某一模式从边界看入的无量纲输入电导为 $y_\mu$，真实电导为 $y_\mu/R$。越过最左边的横向电阻后，
节点处同时接有横向电导 $\mu/R$ 和一个完全相同的无限尾部 $y_\mu/R$。串并联关系给出
\[
  y_\mu=\frac{\mu+y_\mu}{1+\mu+y_\mu},
\]
所以
\[
  y_\mu^2+\mu y_\mu-\mu=0,
  \qquad
  y_1=\frac{\sqrt5-1}{2},
  \qquad
  y_3=\frac{\sqrt{21}-3}{2}.
\]

为说明端口电阻如何组合，令单位模态向量为 $\boldsymbol u_j=\boldsymbol e_j/|\boldsymbol e_j|$，
注入电流为 $I\boldsymbol b$。其模态电流是 $I(\boldsymbol b\cdot\boldsymbol u_j)$，模态电压是
$IR(\boldsymbol b\cdot\boldsymbol u_j)/y_j$。端电压为 $\boldsymbol b\cdot\boldsymbol V$，故
\[
 R_{\rm port}=R\sum_{j=1,3}\frac{(\boldsymbol b\cdot\boldsymbol u_j)^2}{y_j}.
\]
对 $A,C$ 注入与抽出的电流向量正比于 $(1,0,-1)=\boldsymbol e_1$，只激发 $\mu=1$ 模式。该向量
长度平方为 $2$，故
\[
  \boxed{R_{AC}=\frac{2R}{y_1}=(1+\sqrt5)R}.
\]

对 $A,B$，向量 $(1,-1,0)$ 在归一化的 $\boldsymbol e_1,\boldsymbol e_3$ 上的投影平方分别为
$1/2,3/2$。两个正交模式的电压降相加：
\[
  R_{AB}=R\left(\frac{1/2}{y_1}+\frac{3/2}{y_3}\right)
  =\boxed{\frac{4+\sqrt5+\sqrt{21}}4R}.
\]
数值上 $R_{AB}\approx2.70466R$、$R_{AC}\approx3.23607R$。这也符合端点距离直觉：$A,C$ 分处
最上、最下两行，等效电阻较大。
\end{xesolution}



\begin{origexamplebox}{【嵌套网络】}\label{ch21:ex:s37-6-9}
如图所示，框架由同种金属丝制成，单位长度的电阻为 $\rho$。一连串内接等边三角形的数目可认为趋向
无穷，取 $AB$ 边长为 $a$，以后每个三角形的边长依次减小一半。求框架上 $A,B$ 两点间的电阻。
\par\noindent\begin{minipage}{\linewidth}
\nestedTriangleProblemFigure
\diagramcaption{每个新三角形的三个顶点是上一三角形的三边中点，方向逐级交替。}\label{ch21:fig:s37-6-12}
\end{minipage}\par\medskip
\end{origexamplebox}\par\nopagebreak[4]


\begin{xesolution}
任意一层以后仍是一个关于三个端点完全对称的三端网络，所以可把边长为 $a$ 的完整无限网络等效成
三角形连接，设每一条等效支路为
\[
  x(a)=k\rho a,
\]
其中无量纲常数 $k$ 与尺度无关。记 $r=\rho a$。外三角形的每条边被中点分成两个 $r/2$；中点围成的
无限内网边长为 $a/2$，其三角形等效支路就是 $kr/2$。

\par\noindent\begin{minipage}{\linewidth}
\nestedTriangleRecursionFigure
\diagramcaption{解图：把中点以内的无限嵌套部分替换成三条 $kr/2$。外层六段金属丝各为 $r/2$。}
\end{minipage}\par\medskip

在图中令 $A,B$ 电势为 $V,0$。反射对称性与电势反对称性给出 $C$ 及 $AB$ 中点电势为 $V/2$；
令 $CA$ 中点电势为 $u$，则 $CB$ 中点为 $V-u$。在 $CA$ 中点列KCL：
\[
 \frac{u-V}{r/2}+\frac{u-V/2}{r/2}
 +\frac{u-V/2}{kr/2}+\frac{u-(V-u)}{kr/2}=0.
\]
乘 $kr/2$ 后得 $(2k+3)u=3(k+1)V/2$。从 $A$ 流入的总电流为
\[
 I=\frac{V-V/2}{r/2}+\frac{V-u}{r/2}
 =\frac{3(k+2)V}{(2k+3)r}.
\]
所以 $A,B$ 间等效电阻
\[
  R_{AB}=\frac{r(2k+3)}{3(k+2)}.
\]
另一方面，若从完整三端网络的等效三角形看，$A,B$ 之间是一条 $kr$ 与另两条串联的 $2kr$ 并联，
因此
\[
  R_{AB}=kr\parallel 2kr=\frac{2kr}{3}.
\]
令两式相等：
\[
  \frac{2k}{3}=\frac{2k+3}{3(k+2)},
  \qquad 2k^2+2k-3=0.
\]
取正根 $k=(\sqrt7-1)/2$，最终
\[
  \boxed{R_{AB}=\frac{\sqrt7-1}{3}\,\rho a}.
\]
结果小于只保留外三角形时的 $2\rho a/3$，因为不断加入的内接金属丝提供了额外并联通路。
\end{xesolution}



\begin{origexamplebox}{【分形电阻】}\label{ch21:ex:s37-6-10}
如图，考虑大三角形 $ABC$，取每条边的中点并连接，得到第二阶图形，其中含有三个正立小正三角形。
把中间的倒立三角形镂空，便得到一阶“镂垫”。对每一个保留的小三角形重复相同操作，得到更高阶图形。
如此无穷进行，得到谢尔宾斯基三角形，也称谢尔宾斯基镂垫。讨论下列性质：

(1) 最终镂垫的总面积是多少？

(2) 最终镂垫的总边长是多少（内、外边界均计入）？

(3) 若所有边都由单位长度电阻为 $\lambda$ 的电阻丝构成，$AB$ 两点的等效电阻为多少？

(4) 讨论上述结果。
\par\noindent\begin{minipage}{\linewidth}
\sierpinskiProblemFigure
\diagramcaption{前三阶图形由同一中点规则生成；白色倒三角形是被镂空的部分。}\label{ch21:fig:s37-6-13}
\end{minipage}\par\medskip
\end{origexamplebox}\par\nopagebreak[4]


\begin{xesolution}
设最初等边三角形边长为 $a$、面积为 $S_0=\sqrt3a^2/4$。

\textbf{(1) 面积。} 每次操作把一个实心三角形分成四个全等小三角形，只保留其中三个。因此第 $n$ 次
操作后
\[
  S_n=S_0\left(\frac34\right)^n.
\]
令 $n\to\infty$，得到
\[
  \boxed{S_\infty=0}.
\]

\textbf{(2) 总边长。} 第 $n$ 阶有 $3^n$ 个最小正三角形，每个边长 $a/2^n$；这些三角形彼此只在
顶点相接，不共享完整边。把每条几何线段只计一次，总边长为
\[
  L_n=3^n\cdot3\frac{a}{2^n}
  =3a\left(\frac32\right)^n.
\]
所以
\[
  \boxed{L_\infty=\infty}.
\]

\textbf{(3) 等效电阻。} 第 $n$ 阶的最小边长为 $a/2^n$，每段电阻
\[
  r_n=\lambda\frac{a}{2^n}.
\]
若所有最小边的电阻都暂记为 $r$，三角形初阶的端点电阻为 $2r/3$。下面推导增加一阶的比例。
把三个小三角形拼接成一阶网，九条边均为 $r$。令 $A,B$ 电势为 $V,0$，则 $C$ 与 $AB$ 中点为 $V/2$；
$CA,CB$ 中点分别为 $u,V-u$。前一中点连接 $A,C$ 及另外两个中点，KCL给出
\[
 (u-V)+(u-V/2)+(u-V/2)+(u-(V-u))=0,
 \qquad u=3V/5.
\]
所以 $I=(V-V/2)/r+(V-3V/5)/r=9V/(10r)$，一阶端电阻为 $10r/9$。
与初阶 $2r/3$ 相比，比例为 $5/3$。任意更高阶的子网都可等效为对称三端三角形，因而同一比例逐阶适用。
因此
\[
  R_n=\frac23\left(\frac53\right)^n r_n
  =\frac{2\lambda a}{3}\left(\frac56\right)^n.
\]

\par\noindent\begin{minipage}{\linewidth}
\sierpinskiResistanceFigure
\diagramcaption{解图：每一阶由三个缩小到 $1/2$ 的同型三端网络拼成；拓扑使等效电阻乘 $5/3$，长度缩放又使单段电阻乘 $1/2$。}
\end{minipage}\par\medskip

取极限得
\[
  \boxed{R_{AB}=0}.
\]
这是理想无限阶模型的极限，不表示有限粗细的真实金属丝会有严格为零的电阻。

\textbf{(4) 讨论。} 这个分形同时具有“面积趋于零、总边长趋于无穷、端点电阻趋于零”三种看似矛盾
的性质。它们分别测量二维占据、线性边界和电流并联通道，不应互相替代。面积每阶乘 $3/4$；总边长
每阶乘 $3/2$；电阻则是拓扑因子 $5/3$ 与几何缩放 $1/2$ 的乘积，即每阶乘 $5/6$。真实导线有有限
线径、接触电阻和最小结构尺度，递归在某一阶必然停止，因而真实面积、长度和电阻都保持有限且非零。
\end{xesolution}
}

\endgroup

\endgroup
